% Curated book edition: XHIGH (300k 2-part Long Decode) base with cross-references, label compatibility, corrected Part-2 metadata timestamps, and selected proofs/notes integrated from the previous curated Week 7 book edition.
% ==========================================
% Step 4/4 - Final refinement (refinement step 3)
% Output: step4-10-27-2026-tuesday-week7-Algebra_I-offset-final.tex
% Input: step3-10-27-2026-tuesday-week7-Algebra_I-offset-speech_refined.tex
% Model: gemini-4-hydrogen
% Temperature: 0.35
% TopP: 0.9
% TopK: 10
% MaxOutputTokens: 300000
% ThinkingBudget: 4096
% ThinkingLevel: XHIGH
% Processed on: 2026-11-14 01:33:56
% Thought summaries: step4-10-27-2026-tuesday-week7-Algebra_I-offset-final-thoughts.md
% ------------------------------------------
% Token Usage (2 Request(s)):
%   - Total Prompt Tokens : 236'604 (gesamter Prompt, inkl. Cache)
%   - Cached Context      : 0 (aus Google Context-Cache, rabattiert)
%   - Fresh Input Tokens  : 236'604 (neu verrechneter Input)
%   - Output Tokens       : 52'172 (generiertes LaTeX, ohne Thinking)
%   - Thinking Tokens     : 58'764 (Denk-Tokens, als Output verrechnet)
% ==========================================
% ==========================================
% Step 3/4 - Speech refinement (refinement step 2)
% Output: step3-10-27-2026-tuesday-week7-Algebra_I-offset-speech_refined.tex
% Input: step2-10-27-2026-tuesday-week7-Algebra_I-offset-merged.tex
% Model: gemini-4-hydrogen
% Temperature: 0.35
% TopP: 0.9
% TopK: 10
% MaxOutputTokens: 300000
% ThinkingBudget: 4096
% ThinkingLevel: XHIGH
% Processed on: 2026-11-14 01:24:08
% Thought summaries: step3-10-27-2026-tuesday-week7-Algebra_I-offset-speech_refined-thoughts-2.md
% Non-speech guard: all 66 segments outside speech match the input
% ------------------------------------------
% Token Usage (3 Request(s)):
%   - Total Prompt Tokens : 1'045'029 (gesamter Prompt, inkl. Cache)
%   - Cached Context      : 0 (aus Google Context-Cache, rabattiert)
%   - Fresh Input Tokens  : 1'045'029 (neu verrechneter Input)
%   - Output Tokens       : 51'164 (generiertes LaTeX, ohne Thinking)
%   - Thinking Tokens     : 90'030 (Denk-Tokens, als Output verrechnet)
% ==========================================
% ==========================================
% Step 2/4 - Merge & timestamp alignment (refinement step 1)
% Output: step2-10-27-2026-tuesday-week7-Algebra_I-offset-merged.tex
% Input: step1-10-27-2026-tuesday-week7-Algebra_I-all-offset.tex
% Model: gemini-4-hydrogen
% Temperature: 0.4
% TopP: 0.95
% TopK: 40
% MaxOutputTokens: 300000
% ThinkingBudget: 32765
% ThinkingLevel: XHIGH
% Processed on: 2026-11-14 01:02:29
% Thought summaries: step2-10-27-2026-tuesday-week7-Algebra_I-offset-merged-thoughts.md
% ------------------------------------------
% Token Usage (2 Request(s)):
%   - Total Prompt Tokens : 555'341 (gesamter Prompt, inkl. Cache)
%   - Cached Context      : 0 (aus Google Context-Cache, rabattiert)
%   - Fresh Input Tokens  : 555'341 (neu verrechneter Input)
%   - Output Tokens       : 52'349 (generiertes LaTeX, ohne Thinking)
%   - Thinking Tokens     : 28'592 (Denk-Tokens, als Output verrechnet)
% ==========================================
% ==========================================
% Step 1/4 - Transcription, 2 part(s) combined
% Output: step1-10-27-2026-tuesday-week7-Algebra_I-all-offset.tex
% Source video: 10-27-2026-tuesday-week7-Algebra_I-speed-1-compressed.mp4
% Model: gemini-4-hydrogen
% Temperature: 0.36
% TopP: 0.95
% TopK: 40
% MaxOutputTokens: 300000
% ThinkingBudget: 32765
% ThinkingLevel: XHIGH
% Processed on: 2026-11-14 00:48:35
% ------------------------------------------
% Parts:
%   - Part 1: 10-27-2026-tuesday-week7-Algebra_I-part1.mp4 | start 00:00:00 | model gemini-4-hydrogen | processed 2026-11-14 00:34:19 | prompt 1'241'543 | output 25'374 | thinking 53'173
%   - Part 2: 10-27-2026-tuesday-week7-Algebra_I-part2.mp4 | start 00:42:09 | model gemini-4-hydrogen | processed 2026-11-14 00:48:35 | prompt 1'293'978 | output 30'429 | thinking 50'393
% ------------------------------------------
% Token Usage of this run (4 Request(s)):
%   - Total Prompt Tokens : 2'535'521 (gesamter Prompt, inkl. Cache)
%   - Cached Context      : 0 (aus Google Context-Cache, rabattiert)
%   - Fresh Input Tokens  : 2'535'521 (neu verrechneter Input)
%   - Output Tokens       : 55'803 (generiertes LaTeX, ohne Thinking)
%   - Thinking Tokens     : 103'566 (Denk-Tokens, als Output verrechnet)
% ==========================================

% Primary Language: English
% End of the video: 01:28:18

\lecturechapter[7]{Tuesday}{27 Oct}{27 October 2026}{Gauss's Theorem on UFDs and Introduction to Groups}

\begin{overview}[Course Content Overview]
This lecture finishes the proof of Gauss's theorem on polynomial rings over unique factorization domains and transitions to the study of groups:
\begin{itemize}
    \item \textbf{Gauss's Theorem and Classical UFDs:} Reviewing Theorem~B ($R$ is a UFD $\implies R[x]$ is a UFD), classical polynomial ring examples ($\mathbb{Z}[x], \mathbb{Z}[x,y], \dots$, $\mathbb{F}_p[x,y,z,\dots]$, and $\mathbb{Q}[x,y,z]$), the Gaussian integers $\mathbb{Z}[i]$, and formal power series $k[[x]]$ over a field (proving $k[[x]]$ is a PID and classifying all its ideals).
    \item \textbf{Proof of Gauss's Theorem (Case $\deg(p) > 0$):} Reviewing content, primitive polynomials, and Gauss's Lemma, then proving that every positive-degree irreducible $p \in R[x]$ generates a prime ideal via passing to the fraction field $Q(R)[x]$, clearing denominators, and applying primitivity.
    \item \textbf{Reflections on Studying Mathematics:} Discussing the three pillars of learning the subject (definitions, theorems, and examples) alongside philosophical remarks on independent thinking.
    \item \textbf{Group Theory Foundations:} Reviewing the group axioms, uniqueness of the identity and inverses, left/right cancellation, generalized associativity (and Catalan numbers), the inverse map, subgroups, and group homomorphisms.
    \item \textbf{Symmetric Groups and Element Orders:} Defining $S_n$ ($|S_n| = n!$) and cycle decomposition notation (illustrating non-commutativity in $S_3$), proving every element $x$ in a finite group $G$ has finite order $\operatorname{Ord}(x)$, establishing $|\langle x \rangle| = \operatorname{Ord}(x)$, and stating that $\langle x \rangle \subset G$ is a subgroup isomorphic to $(\faktor{\mathbb{Z}}{\operatorname{Ord}(x)\mathbb{Z}}, +, 0)$.
\end{itemize}
\end{overview}

\section{Gauss's Theorem and Classical Examples of UFDs}

\begin{speech}[00:00:00 - 00:00:47]
Okay, so, the last lecture I was almost at the end of this theorem (Recall: \cref{thm:ufd-polynomial-ring-w6wed}), which I will state again. So, this is Carl Friedrich, and the theorem is: $R$ is a UFD implies the polynomial ring is also a UFD \inlinemetanote{writes \qt{CFG} and \qt{$R \text{ is a UFD} \implies R[x] \text{ is a UFD}$}}.

And there were many steps in the logic here, and maybe I take a chance to review some of the things and also explain some of the ideas. So, this is an intricate proof, and as I said, at the end of the last lecture, I came to within $5$ minutes of completing it. But anyway, I'll take some time now.

So, the first thing to notice --- this is not about the proof... This was what was labeled before Theorem B last week (Recall: \cref{thm:ufd-polynomial-ring}) \inlinemetanote{writes \qt{Thm B} inside a box below \qt{CFG}}.
\end{speech}

\begin{content}
\recap{Theorem~B was stated in Week~6 (\cref{thm:ufd-polynomial-ring}) and restated in the previous lecture (\cref{thm:ufd-polynomial-ring-w6wed}), where the bell cut its proof short. The lecturer restates it and completes the proof below.}

\begin{theorem}[Gauss's Theorem (Theorem B)]\label[theorem]{thm:ufd-polynomial-ring-w7tue}
If $R$ is a unique factorization domain (UFD), then the polynomial ring $R[x]$ is also a UFD:
\[
R \text{ is a UFD} \implies R[x] \text{ is a UFD}.
\]
\end{theorem}

\begin{steps}[Historical Context and Scope of the Theorem]
Going beyond the lecture: Gauss proved a form of the key lemma on primitive polynomials (\cref{lem:gauss-lemma}) for integer polynomials in his \emph{Disquisitiones Arithmeticae} (1801, Art.~42); the statement for an arbitrary UFD $R$ is the modern form of his argument. By induction on the number of variables $n$, the theorem extends to multivariate polynomial rings (as in \cref{ex:multivariate-poly-ufd}):
\[
R[x_1, \dots, x_n] \cong (R[x_1, \dots, x_{n-1}])[x_n].
\]
\end{steps}
\end{content}

\subsection{Standard Examples of Unique Factorization Domains}

\begin{speech}[00:00:47 - 00:02:16]
So, the first thing to note, which is not about the proof, it's just about examples \inlinemetanote{writes: Examples} --- that if you know this result, then you can take any UFD you want, like $\mathbb{Z}$ is a UFD, so this result implies if I take the polynomial ring in one variable, or two variables, or any number of variables \inlinemetanote{writes \qt{$\mathbb{Z}$} and \qt{$\implies \mathbb{Z}[x], \mathbb{Z}[x, y], \dots$}} --- you can just keep applying this one after the other \inlinemetanote{points at \cref{thm:ufd-polynomial-ring-w7tue} on the board} --- that these are all UFDs.

And also --- these are kind of the standard classical examples --- if I take the finite field $\mathbb{Z} \bmod p\mathbb{Z}$: every field is a UFD, so that means I take polynomials in finite fields, these, however many variables you want, as long as it's finitely many \inlinemetanote{writes \qt{$\mathbb{F}_p = \faktor{\mathbb{Z}}{p\mathbb{Z}}$} and \qt{$\implies \faktor{\mathbb{Z}}{p\mathbb{Z}}[x, y, z, \dots]$}}. These obey all the factorization rules that you are familiar with for the integers. The idea of: you have a polynomial, there's unique factorization into prime polynomials, and just whatever you remember about the integers --- not every single statement, but the parts about unique factorization, they apply here.

And the other classical examples: you take the fields $\mathbb{Q}, \mathbb{C}, \mathbb{R}$, these are like the classical examples \inlinemetanote{writes \qt{$\mathbb{Q}, \mathbb{C}, \mathbb{R}$} and \qt{$\implies \mathbb{Q}[x, y, z]$}}. And one of the aspects of this course is that, if you see these rings --- there's a huge study of these rings, and this course is only the very beginning of this study ---, but at least you'll know that you can factor things, and the rules of factorization work like in integers.
\end{speech}

\begin{content}
\begin{examples}[Classical Examples of UFDs]
\label[examples]{exm:classical-ufds}
Applying \cref{thm:ufd-polynomial-ring-w7tue} iteratively to standard UFDs yields unique factorization domains in finitely many indeterminates:
\begin{enumerate}[label=\textbf{\arabic*.}]
    \item \textbf{Polynomials over the Integers:} Since $\mathbb{Z}$ is a UFD,
    \[
    \mathbb{Z} \implies \mathbb{Z}[x], \; \mathbb{Z}[x, y], \dots
    \]
    are all UFDs.
    \item \textbf{Polynomials over Finite Fields:} For any prime $p$, the finite field $\mathbb{F}_p = \faktor{\mathbb{Z}}{p\mathbb{Z}}$ is a field (\cref{prop:z-mod-m-domain}) and hence a UFD, so
    \[
    \mathbb{F}_p = \faktor{\mathbb{Z}}{p\mathbb{Z}} \implies \faktor{\mathbb{Z}}{p\mathbb{Z}}[x, y, z, \dots]
    \]
    is a UFD in any finite number of variables.
    \item \textbf{Polynomials over Standard Number Fields:} The fields $\mathbb{Q}, \mathbb{C}, \mathbb{R}$ are UFDs, so polynomial rings such as
    \[
    \mathbb{Q}, \mathbb{C}, \mathbb{R} \implies \mathbb{Q}[x, y, z]
    \]
    (and likewise over $\mathbb{C}$ and $\mathbb{R}$ in any finite number of variables) are UFDs.
\end{enumerate}
\end{examples}

\begin{steps}[Iterative Application of Gauss's Theorem]
For two variables, $R[x, y] \cong (R[x])[y]$. If $R$ is a UFD, \cref{thm:ufd-polynomial-ring-w7tue} guarantees that $R[x]$ is a UFD, and applying \cref{thm:ufd-polynomial-ring-w7tue} a second time with $R[x]$ as the base ring shows that $(R[x])[y]$ is a UFD. By induction, any polynomial ring $R[x_1, \dots, x_n]$ in finitely many variables over a UFD $R$ is again a UFD.
\end{steps}
\end{content}

\begin{aiNote}[Beyond the Lecture: Contrast Between UFDs and PIDs]
% [PERSISTENT, ADDITION]
While the polynomial ring $k[x]$ in one variable over a field $k$ is a principal ideal domain (PID) by division with remainder (\cref{prop:kx-ideals-principal}), the rings $\mathbb{Z}[x]$ (coefficients not in a field) and $k[x, y]$ (two variables) are \emph{not} PIDs: the ideals $(2, x) \subset \mathbb{Z}[x]$ and $(x, y) \subset k[x, y]$ cannot be generated by any single polynomial (\cref{exer:prove-2-x-not-principal}, \cref{ex:non-pids}). Nevertheless, by Gauss's Theorem (\cref{thm:ufd-polynomial-ring-w7tue}), both rings are unique factorization domains.
\end{aiNote}

\subsection{The Gaussian Integers \texorpdfstring{$\mathbb{Z}[i]$}{Z[i]}}

\begin{speech}[00:02:16 - 00:03:11]
Okay, so, you could ask, say: what about other rings? So, it was kind of an example in this class, this ring of Gaussian integers (Recall: \cref{def:gaussian-integers}) \inlinemetanote{writes: $\mathbb{Z}[i]$}. What about this ring? That's not exactly in this context, and it's not... So, it turns out this --- Gaussian integers --- is a PID \inlinemetanote{writes: $\text{is PID}$}. We haven't proven it. You can go look in the book. It was not in the topics in the course, but this is a PID, so it implies that this ring, the Gaussian integers, is also a UFD \inlinemetanote{writes: $\implies \mathbb{Z}[i] \text{ is UFD}$}.

So, the argument it's a PID is not hard; it's just a page in the book. There's a notion of Euclidean domains, and well, you need to know --- so, you have to look in the book for it --- there's a notion of a division algorithm here \inlinemetanote{writes \qt{division alg} below \qt{is PID}}. Then that shows it's a PID, and so, if you're interested, you can read it, you can find it; it's a page in the book basically. But that's also a UFD.
\end{speech}

\begin{content}
\begin{example}[{The Gaussian Integers $\mathbb{Z}[i]$}]
\label[example]{ex:gaussian-integers-ufd}
The ring of Gaussian integers $\mathbb{Z}[i]$ (\cref{def:gaussian-integers}) possesses a division algorithm (it is a Euclidean domain), so (\cref{thm:pid-implies-ufd}):
\[
\mathbb{Z}[i] \text{ is a PID} \implies \mathbb{Z}[i] \text{ is a UFD}.
\]
\end{example}
\end{content}

\begin{aiNote}[Beyond the Lecture: Why \texorpdfstring{$\mathbb{Z}[i]$}{Z[i]} is a Euclidean Domain]
% [PERSISTENT, ADDITION]
The Euclidean norm on $\mathbb{Z}[i]$ is $N(a + bi) := a^2 + b^2$ for $a, b \in \mathbb{Z}$. Given $\alpha, \beta \in \mathbb{Z}[i]$ with $\beta \neq 0$, writing $\alpha/\beta = u + vi \in \mathbb{Q}[i]$ ($u, v \in \mathbb{Q}$) and rounding $u, v$ to the nearest integers $m, n \in \mathbb{Z}$ gives $q := m + ni \in \mathbb{Z}[i]$ and $r := \alpha - q\beta$ with $N(r) \le \bigl(\frac{1}{4} + \frac{1}{4}\bigr) N(\beta) = \frac{1}{2} N(\beta) < N(\beta)$. Hence, $\mathbb{Z}[i]$ is Euclidean, therefore a PID, and therefore a UFD.
\end{aiNote}

\subsection{Formal Power Series over a Field \texorpdfstring{$k[[x]]$}{k[[x]]}}

\begin{speech}[00:03:11 - 00:03:56]
What about rings like... if I take a field, $\mathbb{Q}$ --- let's say $k$: $k$ is a field \inlinemetanote{writes: $k \text{ field}$}. What about the rings --- so, these are polynomial rings \inlinemetanote{points at $\mathbb{Z}[x], \mathbb{Z}[x, y]$} --- what about power series \inlinemetanote{writes: $k[[x]]?$}?

That's not covered here. And it's actually an interesting, complicated topic, but in this particular case it's not complicated. And the answer --- and I hinted at this (Recall: \cref{exer:power-series-pid}) --- is that actually this power series in a field is a PID, so it's also a UFD \inlinemetanote{writes \qt{Ans: $k[[x]] \text{ is PID}$} and \qt{$\implies \text{UFD}$}}.

And I said that I was not going to give a hint about the proof, but I will now basically give the proof. More or less.
\end{speech}

\begin{content}
\begin{proposition}[{Principal Ideal Property of $k[[x]]$}]
\label[proposition]{ex:power-series-pid-ufd}
Let $k$ be a field. Then the ring of formal power series $k[[x]]$ (\cref{def:power-series-ring-formal}) is a principal ideal domain (PID), and consequently a UFD (\cref{thm:pid-implies-ufd}):
\[
k[[x]] \text{ is a PID} \implies k[[x]] \text{ is a UFD}.
\]
(This was posed in Week~6 as \cref{exer:power-series-pid}; the lecturer now sketches the proof.)
\end{proposition}
\end{content}

\begin{proofWrapper}[{Proof of \cref{ex:power-series-pid-ufd}}]

\begin{speech}[00:03:56 - 00:05:31]
So, proof \inlinemetanote{writes: Proof:}. Let $I$ --- I have to prove it's a PID, because we proved earlier that PIDs are UFDs (Recall: \cref{thm:pid-implies-ufd}). So, let $I$ be any ideal \inlinemetanote{writes: Let $I \subset k[[x]]$ ideal}.

Then either $I$ could be $0$ \inlinemetanote{writes: $I = (0)$} --- that's one possibility --- or $I$ is not the zero ideal \inlinemetanote{writes: $I \neq (0)$}. Here, of course, there's nothing to prove if it's the zero ideal.

So, if it's not the zero ideal, choose a power series in here --- choose $f$ in here that has minimal lowest degree \inlinemetanote{writes \qt{$\implies$ choose $f \in k[[x]]$} and \qt{that has min lowest degree} below, later adding \qt{$\in I$} under $k[[x]]$}. Power series don't have highest degrees; you don't think about... you think about lowest degree. So, it's the minimal lowest degree; that means $f$ starts somewhere, like $x^{17} + x^{18} + 3x^{19}$ \inlinemetanote{writes: $f = x^{17} + x^{18} + 3x^{19} + \dots$} --- so, $f$ starts somewhere. And so here, or it could be, you know, there's some element in the field; anyway, something non-zero in the field \inlinemetanote{inserts \qt{$2$} before $x^{17}$}. It starts somewhere, and you find the element here in $I$, so $f$ is in $I$ of lowest degree.

Okay? So, do you accept you can do that? Once it's non-zero, it has some elements, and then you can find the one of lowest degree. All right.
\end{speech}

\begin{content}[Minimal Lowest-Degree Element in an Ideal]
Let $I \subset k[[x]]$ be an arbitrary ideal. We distinguish two cases:
\begin{itemize}
    \item \textbf{Case $I = (0)$:} The zero ideal is generated by the single element $0$, so there is nothing to prove.
    \item \textbf{Case $I \neq (0)$:} Choose a non-zero formal power series $f \in I$ of minimal lowest degree among all non-zero elements of $I$. For instance, $f$ begins at some smallest index with a non-zero coefficient in $k$, such as
    \[
    f = 2x^{17} + x^{18} + 3x^{19} + \dots
    \]
\end{itemize}
\end{content}

\begin{aiNote}[Board Note on the Choice of \texorpdfstring{$f$}{f} and Coefficient \texorpdfstring{$2$}{2}]
% [PERSISTENT]
On the board, the lecturer initially wrote \qt{$\implies$ choose $f \in k[[x]]$} and later added \qt{$\in I$} underneath $k[[x]]$ when clarifying aloud that $f$ is chosen in $I$. Moreover, the example coefficient $2$ stands as an arbitrary non-zero element of $k$ (assuming $\operatorname{char}(k) \neq 2$); in general, the leading lowest-degree term is $a_d x^d$ with $a_d \in k \setminus \{0\}$.
\end{aiNote}

\begin{speech}[00:05:31 - 00:05:44]
Then this is --- now, this is the part for you, exercise \inlinemetanote{writes: Ex: prove $I =$}. I did most of it. Prove $I$ --- what am I going to write here? Yeah?
\end{speech}

\begin{student}[Student Answer]
It's generated by $f$.
\end{student}

\begin{speech}[00:05:46 - 00:06:56]
Yeah, you could say that, but there's a more elegant way to say it. I went to all this trouble writing the minimal degree here. It's generated by the lowest term, in this case $x^{17}$ \inlinemetanote{writes: $(x^{17})$}. It's generated by the lowest term.

So, in general, if you have $f$ is in $I$ of minimal degree, then this is generated by the monomial $x$ to the minimal \inlinemetanote{writes: $f \in I$ min.\ degree $\implies I = (x^m)$} --- that's the one, the monomial where the first one comes, okay? And the reason is that --- well, I mean, we discussed these things, I'll leave that as an exercise for you to think.

This shows that not only is this power series ring a principal ideal domain, because every ideal is generated by one element --- in this case $17$ or some other power ---, but actually I just told you all the ideals in the ring: all the ideals are of the form $x$ to the some number; that's it, those are the only ideals. So, there are very few --- I mean, there's infinitely many, but in some conceptual sense very few ideals.

Okay, those are all PIDs. Sorry --- these rings are all PIDs, so they're all UFDs. Okay, now let's go back to the proof \inlinemetanote{erases the middle and right boards, leaving Theorem B on the upper left}.
\end{speech}

\begin{content}
\begin{exercise}[{Generator of a Non-Zero Ideal in $k[[x]]$}]
\label[exercise]{exer:ideals-power-series}
For $f \in I \setminus \{0\}$ of minimal lowest degree (for instance $f = 2x^{17} + x^{18} + 3x^{19} + \dots$), prove that
\[
I = (x^{17}),
\]
and in general, if $f \in I$ has minimal lowest degree $m$, then
\[
I = (x^m).
\]
\end{exercise}

\begin{proof}[Proof of \cref{exer:ideals-power-series} (beyond the lecture)]
The lecture leaves the proof that $I = (x^m)$ as an exercise. Let $m$ be the minimal lowest degree among non-zero elements of $I$, and write $f = x^m \sum_{j=0}^\infty a_{m+j} x^j = x^m u$ with $a_{m+j} \in k$ and $a_m \in k \setminus \{0\}$. Since $k$ is a field and the constant term $a_m$ of $u$ is non-zero, $u$ is a unit in $k[[x]]$ (\cref{thm:units-power-series-ring}). Therefore, $x^m = u^{-1} f \in I$, which gives $(x^m) \subseteq I$.

Conversely, let $h = \sum_{j=0}^\infty c_j x^j \in I$ be any element. If $h = 0$, then $h \in (x^m)$ trivially. If $h \neq 0$, its lowest degree $n$ satisfies $n \ge m$ by the minimality of $m$. Factoring out $x^m$ yields $h = x^m \sum_{j=n}^\infty c_j x^{j-m} \in (x^m)$. Thus, $I = (x^m)$.
\end{proof}
\end{content}

\end{proofWrapper}

\begin{insight}[{All Ideals of $k[[x]]$}]
As the lecturer emphasizes, this argument does more than show $k[[x]]$ is a PID: it classifies every ideal in the ring. Every non-zero ideal of $k[[x]]$ is of the form $(x^d)$ for some $d \ge 0$, so although there are infinitely many ideals, in a conceptual sense there are very few.
\end{insight}

\begin{hidden}
% timestamp: 00:06:56
% topic: Completed review of Gauss's theorem statement, examples (Z, F_p, Q, C, R, Z[i], k[[x]]), and proof that k[[x]] is a PID
% board_state: thm:ufd-polynomial-ring-w7tue, exm:classical-ufds, ex:gaussian-integers-ufd, ex:power-series-pid-ufd, exer:ideals-power-series
% next_goal: Review content, primitive polynomials, and the two properties from Gauss's Lemma
% open_loops: none
\end{hidden}

\section{Preliminaries: Content, Primitive Polynomials, and Gauss's Lemma}

\begin{speech}[00:06:56 - 00:08:38]
So, the ideas of this proof --- there were several new ideas, which I just comment on, so we're going to use them. If I have $f$ as a polynomial and $f$ is not zero \inlinemetanote{writes: $f \in R[x], \; f \neq 0$} --- we don't like factoring zeros ---, if $f$'s not zero, then we define the content \inlinemetanote{writes \qt{$C(f)$} and \qt{Content of $f$} (in quotes) below}. This is the content of $f$; that's the words for it, and it equals the $\operatorname{GCD}$ of the non-zero coefficients of $f$ \inlinemetanote{writes: $= \operatorname{GCD}(\text{nonzero coeffs of } f)$}. Okay? Since $f$ is not zero, it has at least one coefficient that's non-zero. And that's a $\operatorname{GCD}$, and it's only defined up to a unit \inlinemetanote{adds \qt{$u$} above the equals sign}.

And the ways we think about this is: given any $f$ --- always, you know, non-zero \inlinemetanote{writes: $f \in R[x], \; f \neq 0$} ---, $f$ non-zero, then we can write $f$ is equal to the content --- where we just pull out the $\operatorname{GCD}$ \inlinemetanote{writes: $\implies f = C(f)$} --- and then we multiply by what's left, and this is what's called primitive \inlinemetanote{writes \qt{$\cdot f^*$} and draws an arrow from $f^*$ up to \qt{primitive}}. And the definition of the term primitive was: the content is equal to $1$ up to a unit \inlinemetanote{draws an arrow from $f^*$ down to \qt{$C(f^*) \stackrel{u}{=} 1$}}.

Okay? So, this is review from last time (Recall: \cref{def:content-poly}, \cref{def:primitive-poly}). And the two properties that we use about this, which are really important: the first one is another theorem by Carl Friedrich, that says that if $f$ and $g$ are primitive, then their product is primitive \inlinemetanote{writes \qt{$\bullet \; \text{CFG:}$}, \qt{$\text{if } f, g \text{ are primitive}$}, and \qt{$\implies fg \text{ primitive}$}}. Okay? So, we proved that (Recall: \cref{lem:gauss-lemma}).
\end{speech}

\begin{content}
\recap{Content and primitive polynomials were defined in the previous lecture (\cref{def:content-poly}, \cref{def:primitive-poly}), with the canonical decomposition $f = C(f) f^*$ (\cref{prop:canonical-decomp}). Gauss's Lemma (\cref{prop:gauss-lemma-product} below) was proved there as \cref{lem:gauss-lemma}.}

\begin{definition}[Content and Primitive Polynomial]\label[definition]{def:content-primitive}
Let $R$ be a UFD and let $f \in R[x]$ with $f \neq 0$.
\begin{itemize}
    \item The \newterm{content} of $f$, denoted $C(f)$, is the greatest common divisor of the non-zero coefficients of $f$, defined up to a unit in $R$ (denoted $\stackrel{u}{=}$):
    \[
    C(f) \stackrel{u}{=} \operatorname{GCD}(\text{nonzero coefficients of } f).
    \]
    \item Factoring out the content decomposes any non-zero polynomial $f \in R[x]$ as
    \begin{equation}\label{eq:content-factorization}
    f = C(f) \cdot f^*,
    \end{equation}
    where $f^* \in R[x]$ is \newterm{primitive}, meaning that
    \[
    C(f^*) \stackrel{u}{=} 1.
    \]
\end{itemize}
\end{definition}

\begin{proposition}[Gauss's Lemma: Product of Primitive Polynomials]\label[proposition]{prop:gauss-lemma-product}
Let $R$ be a UFD. If $f, g \in R[x]$ are primitive polynomials, then their product $fg$ is also primitive:
\[
f, g \text{ are primitive} \implies fg \text{ is primitive}.
\]
\end{proposition}
\end{content}

\begin{speech}[00:08:38 - 00:10:17]
And that was a very nice proof actually; it used the fact about polynomial rings over integral domains being integral domains (Recall: \cref{thm:domain-poly-power-series}). And what the content of this is, is if you take something like --- let's take the integers --- $(2x$, I don't know, plus something, $5)$ times $(3x + 17)$ \inlinemetanote{writes: $(2x + 5)(3x + 17)$}.

So, like these two things, these two polynomials --- this is $f$ and this is $g$ \inlinemetanote{labels $2x + 5$ as $f$ and $3x + 17$ as $g$ below} ---, this one's primitive, because $2$ and $5$ don't have any common factors; this one's also primitive, because $3$ and $17$ don't have any common factors. And what our friend Carl Friedrich is asserting here is if I multiply these, it'll also have that property.

So, then to check it, we can multiply. So, we have to do some multiplication here: $6x^2$ plus something, plus $17$ times $5$, $85$, and this is the one where you have to add: I think it's $49$ \inlinemetanote{writes: $= 6x^2 + 49x + 85$}. Is that right? Anyway, you can do some multiplication. It's important to get this right. Looks reasonable.

Okay, so, then if you factor this, you get $2 \times 3$, $7^2$, and this is $5 \times 17$, and you see they have no common factors \inlinemetanote{writes $2 \cdot 3$, $7^2$, and $5 \cdot 17$ beneath $6$, $49$, and $85$}. And that's the assertion he's saying: every single time you do this, that happens. You can't...

And the thing that's tricky --- I mean, we proved it, so I'm not repeating the proof ---, but the thing that's tricky here is that although these things are relatively prime, when you multiply, there's all sorts of addition that happens, and he's saying that doesn't stop his claim from being true. Okay, that's one principle.
\end{speech}

\begin{content}
\begin{example}[{Numerical Verification of Gauss's Lemma in $\mathbb{Z}[x]$}]
\label[example]{ex:gauss-lemma-integers}
Consider the two polynomials in $\mathbb{Z}[x]$:
\[
f = 2x + 5, \qquad g = 3x + 17.
\]
Both $f$ and $g$ are primitive because $\operatorname{GCD}(2, 5) = 1$ and $\operatorname{GCD}(3, 17) = 1$. Multiplying them gives:
\[
(2x + 5)(3x + 17) = 6x^2 + 49x + 85.
\]
\checkline{Checking the middle coefficient and prime factorizations: $2 \cdot 17 + 5 \cdot 3 = 34 + 15 = 49$, and factoring each coefficient yields $6 = 2 \cdot 3$, $49 = 7^2$, and $85 = 5 \cdot 17$, which share no common factor: $\operatorname{GCD}(6, 49, 85) = 1$, confirming that $fg$ is primitive.}
\end{example}
\end{content}

\begin{speech}[00:10:17 - 00:11:19]
The second one was that if I take... \inlinemetanote{starts writing \qt{$a\,b$}, then erases it} If $g$ divides $b$ times $f$, and the assumptions are: $g$ is primitive --- so that means, you know, $g$ is some polynomial, primitive polynomial ---, $b$ is a scalar not equal to $0$, and $f$ is some polynomial \inlinemetanote{writes \qt{$\bullet \; g \mid bf$}, with \qt{primitive} and \qt{$g \in R[x]$} below $g$, \qt{scalar}, \qt{$\in R$}, and \qt{$b \neq 0$} below $b$, and \qt{$f \in R[x]$} below $f$} --- the fact that $g$ is primitive implies $g$ actually divides $f$ \inlinemetanote{writes: $\implies g \mid f$}.

So, those are the two things that we're going to use. I'm not going to prove them again; I proved them last time (Recall: \cref{lem:gauss-lemma}, \cref{clm:primitive-divisibility}). But if you've never seen the content, it's worth playing with; it's kind of useful --- a useful idea. Ah, here it is.
\end{speech}

\begin{content}
\recap{This is \cref{clm:primitive-divisibility} of the previous lecture.}

\begin{proposition}[Divisibility by a Primitive Polynomial]\label[proposition]{prop:primitive-divisibility}
Let $R$ be a UFD, $g, f \in R[x]$, and $b \in R$ with $b \neq 0$. If $g$ is primitive and $g \mid bf$ in $R[x]$, then $g$ divides $f$ in $R[x]$:
\[
\bigl( g \mid bf \text{ in } R[x], \; g \text{ primitive}, \; b \in R \setminus \{0\} \bigr) \implies g \mid f \text{ in } R[x].
\]
\end{proposition}
\end{content}

\begin{aiNote}[Beyond the Lecture: A Second Proof of the Primitive Divisor Lemma]
% [PERSISTENT, ADDITION]
With \cref{prop:gauss-lemma-product}, \cref{prop:primitive-divisibility} also follows directly: we may assume $f \neq 0$, and write $b f = g h$ for some $h \in R[x]$. Taking contents gives $b \, C(f) \stackrel{u}{=} C(b f) = C(g h) \stackrel{u}{=} C(h)$ because $g$ is primitive (\cref{prop:gauss-lemma-product}, applied to $g$ and $h^*$). Hence $b \mid C(h)$ in $R$, so $h = b h'$ with $h' \in R[x]$, and cancelling $b \neq 0$ in the domain $R[x]$ yields $f = g h'$.
\end{aiNote}

\section{Structure of the Proof of Gauss's Theorem}

\begin{speech}[00:11:19 - 00:12:13]
Okay, so, now we come to the structure of the proof. So, it's intricate, and there are many steps we've already done, which I don't repeat, but I'll just show you the structure. So, we want to prove this: $R$ UFD implies $R[x]$ UFD \inlinemetanote{writes: $R \text{ UFD} \implies R[x] \text{ UFD}$}.

For $R[x]$, we have to prove the existence \inlinemetanote{writes: For $R[x]$, we have to prove} --- that's the existence for factorization into irreducibles \inlinemetanote{writes: $\exists \text{ --- Factorization } \to \text{ irr. \checkmark}$} ---, that's the existence part, and I did that all last time. That was some induction on degree.

Then we also have to prove uniqueness; so, that's the uniqueness statement \inlinemetanote{writes: $! \text{ --- Uniqueness}$}.
\end{speech}

\begin{content}[High-Level Structure: Existence and Uniqueness]
To prove \cref{thm:ufd-polynomial-ring-w7tue} ($R \text{ UFD} \implies R[x] \text{ UFD}$), we must establish two properties for $R[x]$:
\begin{itemize}
    \item \textbf{Existence ($\exists$):} Every non-zero, non-unit polynomial in $R[x]$ factors into irreducibles. (Already proved in the preceding lecture by induction on the degree.)
    \item \textbf{Uniqueness ($!$):} Factorization into irreducibles in $R[x]$ is unique up to order and associates.
\end{itemize}
\end{content}

\begin{speech}[00:12:13 - 00:13:52]
And for uniqueness, we use this principle that we proved last time, and this principle is that, if we have $p$ in $R[x]$ is irreducible \inlinemetanote{opens a large square bracket and writes: $p \in R[x] \text{ irr.}$}: if this implies the ideal in $R[x]$ is prime \inlinemetanote{writes: $\implies (p) \subset R[x] \text{ prime}$} --- so, if I can prove this statement \inlinemetanote{closes the square bracket}, then that implies the uniqueness of factorization into irreducibles \inlinemetanote{draws $\Downarrow$ and writes: $! \text{ --- Factorization into irreducibles}$}.

This is a principle we proved last Tuesday (Recall: \cref{clm:step1-uniqueness}), that if I have a ring such that the irreducibles are prime, then there the uniqueness statement holds, if you remember that. So, we're going to use it now. We need to prove this.

So, reminder: $p$ not equal to $0$ or a unit is irreducible if we have: whenever $p = AB$, then at least one of these is a unit --- then $A$ is a unit or $B$ is a unit \inlinemetanote{writes \qt{Reminder:}, \qt{$p \neq 0, \text{ unit}$ is irr.}, \qt{if $p = AB$}, and \qt{then $A \text{ unit}$ or $B \text{ unit}$}}. That's just the definition of... just I remind you.

We have to prove that whenever I have an irreducible person here \inlinemetanote{points at $p \in R[x] \text{ irr.}$}, the ideal generated is prime. That will prove the uniqueness --- we already have existence --- and then we will have proved Gauss's theorem.
\end{speech}

\begin{content}[Reduction of Uniqueness to Primality of Irreducibles]
\recap{The reduction to prime ideals is \cref{clm:step1-uniqueness} (Week~6 Tuesday), and irreducible elements were defined in \cref{def:irreducible}; the existence part ($\exists$) was proved in the previous lecture.}

By \cref{clm:step1-uniqueness} (existence of factorization into irreducibles together with every irreducible generating a prime ideal implies UFD), uniqueness reduces to proving the implication:
\begin{multline*}
\Bigl[\, p \in R[x] \text{ irreducible} \implies (p) \subset R[x] \text{ prime} \,\Bigr] \\
\implies \text{Uniqueness ($!$) of factorization into irreducibles.}
\end{multline*}
\recap{An element $p \neq 0$ which is not a unit is \emph{irreducible} (\cref{def:irreducible}) if whenever $p = AB$, then $A$ is a unit or $B$ is a unit.}
\end{content}

\begin{hidden}
% timestamp: 00:13:52
% topic: Review of content, primitive polynomials, Gauss's Lemma, and structural reduction of Gauss's Theorem
% board_state: def:content-primitive, prop:gauss-lemma-product, ex:gauss-lemma-integers, prop:primitive-divisibility
% next_goal: Split (p) prime into Case (i) deg(p) = 0 and Case (ii) deg(p) > 0, then prove Case (ii)
% open_loops: none
\end{hidden}

\begin{speech}[00:13:52 - 00:15:33]
All right, so, even this I started last time. Our goal here is: we must prove that implication --- $p$ in $R[x]$ irreducible implies the ideal $(p)$ is prime \inlinemetanote{writes \qt{We must prove} and \qt{$p \in R[x] \text{ irr.} \implies (p) \subset R[x] \text{ is prime}$}}. This is the statement we must prove.

And even for this, there were two cases. Case \textbf{(1)} is that the degree of $p$ is $0$ \inlinemetanote{writes: CASE (i) $\deg(p) = 0$}. So, what does it mean if $p$ is a polynomial of degree $0$? This means really $p$ is a scalar; it's in $R$ \inlinemetanote{writes: $p \in R$}. And we discussed this case last time. But this case should certainly not surprise you, because once $p$ is in $R$, we're assuming $R$ is a UFD, so we just factorize in $R$ --- that's basically the summary of that proof. So, we discussed this last time; I'm not going to discuss it again. This is essentially because $R$ is a UFD, we're just assuming \inlinemetanote{writes: $\longleftarrow R \text{ is UFD}$}. It's not a problem.

So, the interesting case --- and that's actually the most interesting case of the whole proof, which is why I try to do it carefully --- is the degree positive one \inlinemetanote{writes: CASE (ii) $\deg(p) > 0$}. Okay, and that's the statement I was in the middle of, or actually almost at the end of last time. But it is an intricate --- I mean, at least for the kind of arguments we've had in this class, it is a slightly intricate one. So, let's do it slowly, and it's pretty actually. As I said, most of the arguments for these things are proofs that prove themselves. This is not one; you have to actually have some ideas.
\end{speech}

\begin{content}[Goal and Case Split by Degree]
We must prove the implication:
\begin{equation}\label{eq:irred-implies-prime}
p \in R[x] \text{ irreducible} \implies (p) \subset R[x] \text{ is prime}.
\end{equation}
This divides into two cases according to the degree of $p$:
\begin{itemize}
    \item \textbf{Case (i):} $\deg(p) = 0$, i.e., $p \in R$. Since $p$ is irreducible in $R[x]$ and constant, $p$ is irreducible in $R$. Because $R$ is a UFD, $p$ is prime in $R$ (\cref{exer:irred-prime-ufd}), which implies $(p) \subset R[x]$ is prime (already completed in the previous lecture).
    \item \textbf{Case (ii):} $\deg(p) > 0$. This is the core case proved below.
\end{itemize}
\end{content}

\begin{aiNote}[Beyond the Lecture: Quotient Polynomial Ring Proof of Case (i)]
% [PERSISTENT, ADDITION]
Reduction of coefficients modulo $(p)$ defines a surjective ring homomorphism $\pi \colon R[x] \to \bigl(\faktor{R}{(p)}\bigr)[x]$ with $\ker(\pi) = (p) \subset R[x]$, so $\faktor{R[x]}{(p)} \cong \bigl(\faktor{R}{(p)}\bigr)[x]$ by the First Isomorphism Theorem (\cref{thm:first-isomorphism-theorem-statement}). Since $p$ is prime in $R$, the quotient $\faktor{R}{(p)}$ is an integral domain (\cref{prop:quotient-integral-domain-prime}), hence $\bigl(\faktor{R}{(p)}\bigr)[x]$ is also an integral domain (\cref{thm:domain-poly-power-series}), and therefore $(p) \subset R[x]$ is a prime ideal.
\end{aiNote}

\section{Proof of Case (ii): Positive-Degree Irreducible Polynomials}

\begin{proofWrapper}[{Proof of Case (ii): \texorpdfstring{$(p)$}{(p)} is Prime for \texorpdfstring{$\deg(p) > 0$}{deg(p) > 0}}]

\subsection{Setup and Minimal-Degree Element in \texorpdfstring{$I = (f, p)$}{I = (f, p)}}

\begin{speech}[00:15:33 - 00:16:43]
Okay, so, what do we have to do? Let's write the assumptions \inlinemetanote{writes: Assumptions:}. So, one is that $p \in R[x]$ is irreducible, and the degree of $p$ is positive \inlinemetanote{writes: $\bullet \; p \in R[x] \text{ is irr., } \deg(p) > 0$}. That's where we are in the logic.

We want to prove that the ideal is prime, so then we can assume $fg$ is in the ideal, and we assume $f$ is not in the ideal \inlinemetanote{writes \qt{$\bullet \; fg \in (p)$} and \qt{$\bullet \; f \notin (p)$}}. All right? And we need to show --- this is our goal \inlinemetanote{draws an implication arrow labeled \qt{Goal} to the right} --- is to show that then $g$ is in the ideal \inlinemetanote{writes: $g \in (p)$}.

All right, so if we can do this, it says: these are our assumptions, and this last part is it being prime: take any two where the first one's not, then the second one has to be. So, that's what we have to do. Okay. That's the logical thing we have to do. Is that clear?
\end{speech}

\begin{content}[Assumptions and Goal]
To prove that $(p) \subset R[x]$ is a prime ideal when $\deg(p) > 0$, we start from the following assumptions:
\begin{itemize}
    \item $p \in R[x]$ is irreducible and $\deg(p) > 0$,
    \item $fg \in (p)$,
    \item $f \notin (p)$.
\end{itemize}
\textbf{Goal:} Prove that $g \in (p)$.
\end{content}

\begin{speech}[00:16:43 - 00:18:50]
All right, so the proof of this is: the first idea is we construct this ideal generated by $f$ and $p$, and it's an ideal in this polynomial ring \inlinemetanote{writes \qt{Proof:} and \qt{$I := (f, p) \subset R[x]$}}. So, just to be clear here: everything here --- $p, f, g$ --- are polynomials \inlinemetanote{writes inside a box below the assumptions: $p, f, g \in R[x]$}. All right.

And so I can take the ideal here. Now I select this element --- and this is an idea that we've used many times: if I have this ideal in a ring, like in this polynomial ring, it's a very good idea often to select the element of lowest degree. So, now we're here \inlinemetanote{writes: Select $m \in I$ so that $m \neq 0$ is of lowest degree among all non-zero elements of $I$}.

Okay? Again, you know, this ideal has some people in it. In particular, $p$ is irreducible, so it's not zero, so this ideal is not zero. So, it has some non-zero elements; we choose one of lowest degree, which we're allowed to. It's not unique, but we're allowed to choose an element of lowest degree. Are you happy with that?

And now, the subgoal here is that we will show that $m$ is a constant: $m$ is in $R$ \inlinemetanote{writes inside square brackets: $[\text{We will show } m \in R]$}. In other words, $m$ is of degree $0$. We will show that; that's going to be the goal of the next part of the argument.
\end{speech}

\begin{content}[{Constructing $I = (f, p)$ and Selecting $m$ of Lowest Degree}]
Here $p, f, g \in R[x]$. Define the ideal generated by $f$ and $p$:
\begin{equation}\label{eq:ideal-fp}
I := (f, p) \subset R[x].
\end{equation}
Since $p \in I$ is irreducible, $p \neq 0$, so $I \neq (0)$. By the well-ordering of $\mathbb{N}_0$, we may select a non-zero element $m \in I$ of lowest degree among all non-zero elements of $I$:
\begin{equation}\label{eq:minimal-m}
m \in I \setminus \{0\} \quad \text{such that} \quad \deg(m) \le \deg(h) \quad \forall h \in I \setminus \{0\}.
\end{equation}
\textbf{Subgoal:} We will show that $m \in R$ (i.e., $\deg(m) = 0$).
\end{content}

\begin{hidden}
% timestamp: 00:18:50
% topic: Proof of Case (ii): assumptions, ideal I = (f, p), selection of minimal-degree m in I, subgoal m in R
% board_state: eq:irred-implies-prime, eq:ideal-fp, eq:minimal-m
% next_goal: Pass to fraction field Q(R) to apply division algorithm in Q[x] and clear denominators
% open_loops: none
\end{hidden}

\subsection{Division Algorithm in \texorpdfstring{$Q[x]$}{Q[x]} and Clearing Denominators}

\begin{speech}[00:18:50 - 00:20:57]
Now, what we normally do in the case when we have a field is we use the division algorithm (Recall: \cref{thm:poly-long-division}). So, we'd very much like to use the division algorithm, and this is one of the ideas in this proof, and it's part of the reason why the notion of the fraction field was introduced. So, we would like to use the division algorithm \inlinemetanote{writes: We would like to use the division alg. in $R[x]$}.

But we can't, because this is not a field. So, the solution here --- and this is, as I said, one of the ideas in the proof --- is that... So, the problem is that $R$ need not be a field \inlinemetanote{writes: But the problem is that $R$ need not be a field}. And for the division algorithm you need a field because you're dividing.

So, the idea here --- and this idea is to use the field of fractions \inlinemetanote{writes: Idea:}. And as I said, it's kind of why we introduced it. So, let $Q(R)$ be the field of fractions (Recall: \cref{def:field-of-fractions}) \inlinemetanote{writes: Let $Q(R)$ be the field of fractions of $R$}. So, remember, $R$ sits inside --- I'm just going to call it $Q$ --- sits inside $Q$ \inlinemetanote{writes: $R \subset Q$}, and this is the field of fractions. Every element of $Q$ --- an element of $Q$ is some kind of $r$ over $b$ where $b$ is not zero \inlinemetanote{writes: $\ni \frac{r}{b}, \; b \neq 0$}.

All right? So, all of these things we can consider --- we can view, maybe is a better word \inlinemetanote{writes: We can view} --- all of these people, $p, f, g$, to be inside this polynomial ring over this field \inlinemetanote{writes: $p, f, g \in Q[x]$}, because they live inside $R$, so we just view them inside the field. Okay? And the advantage for this is we can run the division algorithm in the field. That's the whole purpose of this idea. \inlinemetanote{video cuts as the filled board is pushed up to expose a clean lower board}
\end{speech}

\begin{content}[{Passing to the Field of Fractions $Q(R)$}]
We would like to apply the division algorithm in $R[x]$, but $R$ need not be a field. Instead, let $Q := Q(R)$ be the field of fractions of the integral domain $R$ (\cref{def:field-of-fractions}):
\[
R \subset Q \ni \frac{r}{b} \quad (r, b \in R, \; b \neq 0).
\]
Since $R \subset Q$, we may view $p, f, g \in R[x]$ inside the polynomial ring $Q[x]$, where the division algorithm (\cref{thm:poly-long-division}) is valid because $Q$ is a field.
\end{content}

\begin{speech}[00:20:58 - 00:23:05]
So, what does it mean we can run the division algorithm? This means that we can take $f$ and $m$ --- I should put $m$ also there in that list, sorry, just so I'm complete \inlinemetanote{walks to the upper right board and inserts \qt{$m,$} before \qt{$p, f, g \in Q[x]$}}. Also, we've already selected $m$ too, so $m$ is on this list. Sorry, what's going on? Wrong one.

We want to apply the division algorithm. We have $f$: I want to divide --- division algorithm in here, for division by $m$ \inlinemetanote{writes: $f \longrightarrow \text{division alg in } Q[x] \text{, div by } m$}. Okay? So, the consequence of that is I can write $f$ is equal to the quotient times $m$ plus some remainder \inlinemetanote{writes: $f = \hat{q} \cdot m + \hat{r}$}. And these hats signify that these elements may be in this ring \inlinemetanote{writes: $\hat{q}, \hat{r} \in Q[x]$}. Right? Because $f$ and $m$ are also in $Q[x]$, but these things we know are actually in the smaller place --- that they're here \inlinemetanote{marks $f$ and $m$ as belonging to $R[x]$}.

Okay, but we can't apply the division algorithm in this $R$ --- polynomials in $R$ ---, but we can with polynomials with coefficients in $Q$. So, we apply the division algorithm and we get this quotient and remainder, but they live here \inlinemetanote{points at $\hat{q}, \hat{r} \in Q[x]$}. Is that clear?

And the condition that comes out of the division algorithm is that either $\hat{r}$ is zero or the degree of $\hat{r}$ is strictly less than the degree of $m$ \inlinemetanote{writes: $\text{either } \hat{r} = 0 \text{ or } \deg(\hat{r}) < \deg(m)$}, okay? Are you happy with that?
\end{speech}

\begin{content}[{Division of $f$ by $m$ in $Q[x]$}]
Viewing $m, p, f, g \in Q[x]$ and dividing $f \in R[x] \subset Q[x]$ by $m \in R[x] \setminus \{0\} \subset Q[x]$ via the division algorithm in $Q[x]$ yields polynomials $\hat{q}, \hat{r} \in Q[x]$ such that
\begin{equation}\label{eq:div-f-qhat}
f = \hat{q} \cdot m + \hat{r}, \qquad \text{with either } \hat{r} = 0 \text{ or } \deg(\hat{r}) < \deg(m).
\end{equation}
\end{content}

\begin{speech}[00:23:05 - 00:25:01]
And now comes the next step: this step which is called \qt{clear the denominators} \inlinemetanote{writes: Clear the denoms}. Which is: we don't want to work in polynomials with coefficients in this field. Somehow the theorem is about polynomials with coefficients in $R$, so we'd like to get back to polynomials with coefficients in $R$.

And the only problem here --- so, this is a polynomial with coefficients in $Q[x]$ (i.e., $Q$) ---, so this has some, you know, some $c_6 x^6 + c_5 x^5 + \dots + c_0$ \inlinemetanote{writes: $\hat{q} = c_6 x^6 + c_5 x^5 + \dots + c_0$}. Now, all of these coefficients are in the field of fractions of $R$ \inlinemetanote{writes: $c_i \in Q(R)$}. So, I can write them instead as elements in $R$: $\frac{a_6}{b_6}$, $\frac{a_5}{b_5}$ --- this is kind of explicit --- $\frac{a_0}{b_0}$, where all of these $a_i, b_i$ are in $R$ and moreover the denominators are not zero \inlinemetanote{writes \qt{$\frac{a_6}{b_6}, \frac{a_5}{b_5}, \dots, \frac{a_0}{b_0}$} above $c_6, c_5, \dots, c_0$ and writes \qt{$a_i, b_i \in R, \; b_i \neq 0$}}, right? So, I can just write it like this.

So, that means that if I multiply by all these denominators, I can have $b$ times $\hat{q}$ is equal to $a_6' x^6 + a_5' x^5 + \dots + a_0'$, where this $b$ is the product of the $b_i$ \inlinemetanote{writes \qt{$b \hat{q} = a_6' x^6 + a_5' x^5 + \dots + a_0'$} and \qt{$b = \prod b_i$}} --- that's clearing the denominators ---, and then these are all in the ring also \inlinemetanote{writes: $a_i' \in R$}.

Okay? That's clearing the denominators. I did this really explicitly, but you know how to clear denominators; we just clear the denominators. And the price of clearing the denominators --- what's good about clearing the denominators is you land in this ring; the price is you have to multiply by $b$.
\end{speech}

\begin{content}[{Clearing Denominators for Polynomials in $Q[x]$}]
Any polynomial $\hat{q} \in Q[x]$, illustrated for degree $6$ by
\[
\hat{q} = c_6 x^6 + c_5 x^5 + \dots + c_0 = \frac{a_6}{b_6} x^6 + \frac{a_5}{b_5} x^5 + \dots + \frac{a_0}{b_0},
\]
has coefficients $c_i = \frac{a_i}{b_i} \in Q(R)$ with $a_i, b_i \in R$ and $b_i \neq 0$ (for $i = 0, \dots, 6$). Defining $b := \prod_{i=0}^{6} b_i \in R \setminus \{0\}$ (non-zero because $R$ is an integral domain and each $b_i \neq 0$, and analogously $b := \prod_{i=0}^{d} b_i$ for general degree $d$) clears all denominators:
\[
b \hat{q} = a_6' x^6 + a_5' x^5 + \dots + a_0' \in R[x], \qquad \text{where } a_i' \in R \text{ for } i = 0, \dots, 6.
\]
\end{content}

\begin{hidden}
% timestamp: 00:25:01
% topic: Division of f by m in Q[x] and explicit method of clearing denominators
% board_state: eq:div-f-qhat
% next_goal: Apply denominator clearing to f = qhat*m + rhat, deduce r' = 0, and conclude m* | f
% open_loops: none
\end{hidden}

\begin{speech}[00:25:01 - 00:26:40]
So, I can just do that here. Take here: I have $f = \hat{q} m + \hat{r}$ \inlinemetanote{writes: $f = \hat{q} m + \hat{r}$}. So, there exists some $b'$ in $R$, $b'$ not equal to $0$, which clears the denominators of this \inlinemetanote{writes: $\exists\, b' \in R, \; b' \neq 0$}. That means I multiply by $b'$: $b' f$, then I get a new $q'$ and a new $r'$ \inlinemetanote{writes: $b' f = q' m + r'$}. And now $b'$ is in $R$, as I said not equal to $0$, and $q'$ and $r'$ are polynomials with coefficients here \inlinemetanote{writes: $0 \neq b' \in R, \; q', r' \in R[x]$}, okay?

I'm saying it very explicitly, but all I do is I take this, there might be some denominators, I clear them --- which, you know, you've been clearing denominators for a long time by now, I assume. Okay. So, are you happy with this?

And now I look at this $r'$. So, $r'$ --- when I clear denominators, I don't change the degree. So, we have the same thing: $r'$ is equal to $0$ --- because when you clear the denominators, if it were $0$, clearing the denominators will keep it $0$ ---, or, so either $r'$ is $0$ or the degree of $r'$ is less than the degree of $m$, because multiplication by a scalar does not change the degree \inlinemetanote{writes \qt{$r' = 0$} and \qt{or $\deg(r') < \deg(m)$} below $r'$}. We took $\hat{r}$ and we cleared the denominators; it does not change the degree.
\end{speech}

\begin{content}[{Clearing Denominators in $f = \hat{q} m + \hat{r}$}]
Multiplying \eqref{eq:div-f-qhat} by a common denominator $b' \in R \setminus \{0\}$ of all coefficients of $\hat{q}$ and $\hat{r}$ yields:
\begin{equation}\label{eq:cleared-f}
b' f = q' m + r', \qquad \text{with } 0 \neq b' \in R \text{ and } q' := b' \hat{q}, \; r' := b' \hat{r} \in R[x].
\end{equation}
Since $R$ is an integral domain and $b' \neq 0$, scalar multiplication by $b'$ preserves zero and degree:
\[
r' = 0 \quad \text{or} \quad \deg(r') < \deg(m).
\]
\end{content}

\begin{speech}[00:26:40 - 00:28:07]
All right, but when I look at this equation, I notice that $m$ is in $I$ and $f$ is in $I$ \inlinemetanote{writes: $m \in I, \; f \in I$}, because $m$ was the lowest element in $I$ and $f$ is born in $I$. So, this implies that $r'$ is in $I$ \inlinemetanote{writes: $\implies r' \in I$}. Right?

But $r'$ cannot be lower than $m$ in $I$, because $m$ is an element of the lowest degree in $I$. So, this possibility is impossible \inlinemetanote{crosses out \qt{$\deg(r') < \deg(m)$}}. So, it must be that $r'$ is $0$ \inlinemetanote{writes: $\implies r' = 0$}.

Which is kind of what we were expecting; that was the argument for a field if you use that. But this is the way to come to this conclusion even when $R$ is not a field, by going to the field. Okay, so, the conclusion here is that we have $b' f = q' m$ \inlinemetanote{writes: $b' f = q' m$}.

All right, and then if we write this in terms of contents, we have $b' f = q' \, C(m) \, m^*$ \inlinemetanote{writes: $b' f = q' \, C(m) \, m^*$}, so the conclusion is that $m^*$ divides $f$ \inlinemetanote{writes inside a box: $m^* \mid f$}. This uses one of our principles (Recall: \cref{prop:primitive-divisibility}), because $m^*$ is primitive \inlinemetanote{writes \qt{prim} below $m^*$}: the primitive is dividing this product, so it has to divide one of them (namely $f$). All right? So, this is one of our conclusions. Okay.
\end{speech}

\begin{content}[{Vanishing of $r'$ and Proof that $m^* \mid f$}]
Since $m \in I$ and $f \in I$, equation \eqref{eq:cleared-f} gives
\[
r' = b' f - q' m \in I.
\]
By \eqref{eq:minimal-m}, no non-zero element of $I$ can have degree strictly smaller than $\deg(m)$. Thus, the alternative $\deg(r') < \deg(m)$ is impossible, forcing $r' = 0$. Substituting $r' = 0$ into \eqref{eq:cleared-f} and factoring $m = C(m) m^*$ via \eqref{eq:content-factorization} gives:
\begin{equation}\label{eq:bprime-f}
b' f = q' m = q' C(m) m^*.
\end{equation}
Hence, $m^* \mid b' f$ in $R[x]$. Because $m^*$ is primitive and $b' \in R \setminus \{0\}$, \cref{prop:primitive-divisibility} implies:
\begin{equation}\label{eq:mstar-div-f}
m^* \mid f \quad \text{in } R[x].
\end{equation}

\begin{steps}[Why $m^* \mid b'f$ Implies $m^* \mid f$]
From $b' f = \bigl(q' C(m)\bigr) \cdot m^*$, we see directly that $m^*$ divides the product $b' f$ in $R[x]$ with quotient $q' C(m) \in R[x]$. Since $m^*$ is the primitive part of $m \neq 0$, it is primitive, and $b'$ is a non-zero scalar in $R$. Applying \cref{prop:primitive-divisibility} directly cancels the scalar $b'$ and yields $m^* \mid f$.
\end{steps}
\end{content}

\begin{speech}[00:28:07 - 00:29:05]
So, this proof is now going at about half the speed as I was going last time. Which I don't know is good or bad. \inlinemetanote{erases the lower board silently}

You can also just watch the video at half the speed, I guess --- or twice the speed, or whatever you want. But I'm telling you more details of the argument.
\end{speech}

\subsection{Repeating for \texorpdfstring{$p$}{p} and Using Irreducibility to Deduce \texorpdfstring{$m \in R$}{m in R}}

\begin{speech}[00:29:05 - 00:30:33]
All right, so, we did this whole board with the division algorithm; we divided $f$ by $m$, so do the whole board again: repeat for $p$ divided by $m$ \inlinemetanote{writes: Repeat for $p$ divided by $m$:}. So, what happened there?

We take now $p$ here, and now we have some different quotients --- they're double hats because I already used one hat \inlinemetanote{writes: $p = \hat{\hat{q}}\, m + \hat{\hat{r}}$} ---, and again, these two are in the polynomial ring with coefficients in the field \inlinemetanote{writes: $\hat{\hat{q}}, \hat{\hat{r}} \in Q[x]$}, just as they were there \inlinemetanote{points at the upper board}. We just do the same division algorithm.

And now we clear, just as before --- clear denominators \inlinemetanote{writes: Clear denoms} ---, and the conclusion is we get $b$ double hat --- double prime $p$ is equal to $q'' m + r''$, where now all the same conclusions: $0$ is not equal to $b''$ in $R$, and $q'', r''$ are here \inlinemetanote{writes \qt{$b'' p = q'' m + r''$} and \qt{$0 \neq b'' \in R, \; q'', r'' \in R[x]$}}, okay?

So, here I just do the same thing I did upstairs; I just do it with $p$ instead of $f$, and I reach the same conclusion, of course with different possible coefficients.
\end{speech}

\begin{content}[{Division of $p$ by $m$ and Clearing Denominators}]
Repeating the exact same division and denominator-clearing procedure for $p$ divided by $m$ in $Q[x]$ gives
\[
p = \hat{\hat{q}}\, m + \hat{\hat{r}} \qquad \bigl(\hat{\hat{q}}, \hat{\hat{r}} \in Q[x], \; \hat{\hat{r}} = 0 \text{ or } \deg(\hat{\hat{r}}) < \deg(m)\bigr).
\]
Clearing denominators via a common non-zero scalar $b'' \in R \setminus \{0\}$ yields:
\begin{equation}\label{eq:cleared-p}
b'' p = q'' m + r'',
\end{equation}
with $0 \neq b'' \in R$, $q'', r'' \in R[x]$, and $r'' = 0$ or $\deg(r'') < \deg(m)$.
\end{content}

\begin{hidden}
% timestamp: 00:30:33
% topic: Repeated division algorithm and clearing denominators for p
% board_state: eq:bprime-f, eq:mstar-div-f, eq:cleared-p
% next_goal: Deduce r'' = 0, m* | p, and use irreducibility of p to show m in R
% open_loops: none
\end{hidden}

\begin{speech}[00:30:33 - 00:31:57]
But then again, the next step was we notice that $m$ is in $I$ and $p$ is in $I$ \inlinemetanote{writes: $m \in I, \; p \in I$}, because $p$ is also born inside $I$ --- remind you, $I$ is $(f, p)$ \inlinemetanote{writes inside a box above: $I = (f, p)$}. And so the same conclusion holds, that $r''$ is in $I$, and as a consequence, $r''$ must be zero by the same logic \inlinemetanote{writes \qt{$\implies r'' \in I$} and \qt{$\implies r'' = 0$}}.

So, we get the equation --- conclude \inlinemetanote{writes: Conclude}, just as above; in fact, this whole board is just writing that whole argument again --- that $b'' p$ is equal to $q'' m$ \inlinemetanote{writes: $b'' p = q'' m$}. Okay?

We conclude here that this $m^*$ has to divide $p$ \inlinemetanote{writes inside a box: $m^* \mid p$}, for the same reason that this $m^*$ divides $f$ \inlinemetanote{points at \qt{$m^* \mid f$} on the upper board}, because $m$... Well, we can write it further the next step if you want, but it's the same thing: $q'' \, C(m) \, m^*$; this is primitive, it divides the product, so it must divide $p$ \inlinemetanote{writes \qt{$= q'' \, C(m) \, m^*$} after $b'' p = q'' m$ and writes \qt{prim} below $m^*$}, all right? So, it's just the same. It's not surprising the conclusion's the same, because the steps are the same.
\end{speech}

\begin{content}[{Vanishing of $r''$ and Proof that $m^* \mid p$}]
Since $m \in I$ and $p \in I = (f, p)$, equation \eqref{eq:cleared-p} implies $r'' = b'' p - q'' m \in I$. By minimality of $\deg(m)$ in \eqref{eq:minimal-m}, we must have $r'' = 0$. Therefore:
\begin{equation}\label{eq:bdoubleprime-p}
b'' p = q'' m = q'' C(m) m^*.
\end{equation}
Because $m^*$ is primitive and $b'' \in R \setminus \{0\}$, \cref{prop:primitive-divisibility} gives:
\begin{equation}\label{eq:mstar-div-p}
m^* \mid p \quad \text{in } R[x].
\end{equation}
\end{content}

\begin{speech}[00:31:57 - 00:33:18]
All right, so, what does this mean? This means $p$ is equal to $m^*$ times something \inlinemetanote{writes: $p = m^* \cdot B$}. That's what \qt{$m^*$ divides $p$} is: it means $p$ is equal to $m^*$ times something.

So, now we have to use the irreducibility, all right? Since $p$ is irreducible, this implies one of these: either $m^*$ is a unit or $B$ is a unit \inlinemetanote{writes \qt{$p \text{ irr.}$}, \qt{$\implies m^* \text{ unit} \in R[x]$}, and \qt{or $B \text{ unit} \in R[x]$}}. Since it's irreducible --- we have to use it's irreducible now ---, one of these is true.

And the claim is: $B$, this is impossible \inlinemetanote{crosses out \qt{$B$}}. Because if $B$ were a unit, then $p$ would be $m^*$ up to a unit. But $m^*$ divides $f$, so if this were true, then $p$ would divide $f$ \inlinemetanote{writes below: $\implies p \mid f$}. That would mean $f$ is in the ideal generated by $p$ \inlinemetanote{writes: $\implies f \in (p)$}. But our hypothesis is exactly: $f$ is not in the ideal generated by $p$. This contradicts that hypothesis \inlinemetanote{writes \qt{CONTRA} in red and crosses out that branch}. Okay?
\end{speech}

\begin{content}[{Using Irreducibility of $p$ to Rule Out $B$ Being a Unit}]
From \eqref{eq:mstar-div-p}, there exists $B \in R[x]$ such that
\begin{equation}\label{eq:p-mstar-B}
p = m^* \cdot B.
\end{equation}
Since $p$ is irreducible in $R[x]$, one of the two factors must be a unit in $R[x]$:
\[
p \text{ irreducible} \implies m^* \text{ is a unit in } R[x] \quad \text{or} \quad B \text{ is a unit in } R[x].
\]
We rule out the second case by contradiction:
\begin{itemize}
    \item Suppose $B$ were a unit in $R[x]$. Then $m^* = p B^{-1}$, so $p \mid m^*$.
    \item Combining this with $m^* \mid f$ from \eqref{eq:mstar-div-f} yields $p \mid f$, i.e., $f \in (p)$.
    \item This contradicts our assumption that $f \notin (p)$.
\end{itemize}
\end{content}

\begin{speech}[00:33:18 - 00:34:35]
So, here is where we're using irreducibility --- here we start using the hypotheses. Irreducible means we have this dichotomy, and then when we go on the second branch, $B$ is a unit, we get a contradiction: it contradicts that --- so, one of the hypotheses, if you remember, assumption on that board \inlinemetanote{points at the upper board} was that $f$ is not in $(p)$ \inlinemetanote{writes \qt{(Ass $f \notin (p)$)} in red}; that was one of our assumptions. But if you go on this branch, we get $f$ is in $(p)$.

It has to be that $m^*$ is a unit \inlinemetanote{points at \qt{$m^* \text{ unit} \in R[x]$}}, okay? How can $m^*$ be a unit? $m^*$ is a polynomial. The only way it can be a unit is if it's a constant. So, this implies --- that's the conclusion we were seeking: conclusion, $m^*$ is in $R$ \inlinemetanote{writes: Conclude $m^* \in R$}. The only way you can be a unit in a polynomial ring over an integral domain is if you're a constant (Recall: \cref{thm:units-polynomial-ring-integral-domain}). Right?

And this implies that $m$ was actually in $R$ too, because they differ just by the content \inlinemetanote{writes: $\implies m \in R$}. And that's what I said the goal here is. So, after all this work, we selected this $m$ to be an element in the ideal of lowest degree, and we concluded that it has to be in $R$. So, we're almost done now.
\end{speech}

\begin{content}[{Concluding $m^* \in R$ and $m \in R$}]
Since $B$ cannot be a unit, $m^*$ must be a unit in $R[x]$. Because the only units of a polynomial ring over an integral domain $R$ are units of the base ring $R$ (\cref{thm:units-polynomial-ring-integral-domain}), it follows that $m^* \in R$ (in fact $m^* \in \mathcal{U}_R$). Since $m = C(m) \cdot m^*$ with $C(m) \in R$, we conclude:
\begin{equation}\label{eq:m-in-R}
m \in R \setminus \{0\}.
\end{equation}
\end{content}

\begin{speech}[00:34:35 - 00:35:26]
\inlinemetanote{erases the lower board} I think I tried to give this whole argument in like $10$ minutes last time, roughly speaking. So, we are very close to the end of the argument.
\end{speech}

\subsection{Multiplying by \texorpdfstring{$g$}{g}, Primitivity of \texorpdfstring{$p$}{p}, and Conclusion}

\begin{speech}[00:35:26 - 00:36:27]
Okay, so: we have $m$ is in the ideal \inlinemetanote{writes: $m \in I = (f, p)$} --- this is the ideal generated by $f$ and $p$. So, this means we have $m = Af + Bp$ \inlinemetanote{writes: $m = Af + Bp$}, where everything here --- $m, A, f, B, p$ --- are in the ring \inlinemetanote{writes: $m, A, f, B, p \in R[x]$}. Okay? So, this is just because it's in the ideal generated by them.

But now we know something more: we also know that this is in $R$, this is a scalar \inlinemetanote{writes \qt{$0 \neq m \in R$} to the left of $m$}. This $m$ is in $R$; it's a non-zero element of $R$. Okay?

So, we have to get $g$ in the picture. The whole goal is to prove $g$ is in the ideal generated by $p$, and we haven't seen $g$ at all. So, how are we going to get $g$ in this picture? This is the trick we've used many, many times. Yeah?
\end{speech}

\begin{student}[Student Answer]
Multiply by $g$.
\end{student}

\begin{speech}[00:36:28 - 00:37:11]
Multiply by $g$, exactly. We use this many times. Multiply by $g$ \inlinemetanote{writes: $mg = Afg + Bgp$}. Okay? And this is not equal to $0$ \inlinemetanote{writes \qt{$0 \neq$} below $m$}. Okay.

Now, we know this is in the ideal generated by $p$ \inlinemetanote{draws a bracket under $Afg$ labeled \qt{$(p)$}}. That was one of our assumptions: $fg$ is in the ideal generated by $p$. So, this term is in the ideal generated by $p$. This term is in the ideal generated by $p$ \inlinemetanote{circles $Bgp$}. So, this means that $p$ divides $m$ times $g$; $m$ is a non-zero scalar \inlinemetanote{writes $p \mid m \cdot g$ with $\neq 0$ below $m$}. Okay?
\end{speech}

\begin{content}[{Writing $m = Af + Bp$ and Multiplying by $g$}]
Since $m \in I = (f, p)$, there exist polynomials $A, B \in R[x]$ such that
\begin{equation}\label{eq:m-linear-combo}
m = Af + Bp \qquad (m, A, f, B, p \in R[x]),
\end{equation}
and by \eqref{eq:m-in-R}, $m$ is a non-zero scalar ($0 \neq m \in R$). Multiplying \eqref{eq:m-linear-combo} by $g$ yields:
\begin{equation}\label{eq:mg-combo}
mg = Afg + Bgp.
\end{equation}
Since $fg \in (p)$ by assumption, we have $Afg \in (p)$, and clearly $Bgp \in (p)$. Therefore, $mg \in (p)$, which means:
\begin{equation}\label{eq:p-div-mg}
p \mid mg \quad \text{in } R[x], \qquad \text{with } m \in R \setminus \{0\}.
\end{equation}
\end{content}

\begin{hidden}
% timestamp: 00:37:11
% topic: Linear combination m = Af + Bp, multiplying by g, deducing p | mg
% board_state: eq:m-in-R, eq:mg-combo, eq:p-div-mg
% next_goal: Prove p irr. and deg(p) > 0 => p is primitive, and conclude p | g
% open_loops: none
\end{hidden}

\begin{speech}[00:37:11 - 00:39:17]
So, now the last thing is: what is the assumption on $p$? $p$ is irreducible \inlinemetanote{writes: $p \text{ is irr.}$}. If we knew $p$ were primitive, we'd be done \inlinemetanote{writes: If $p$ is primitive:}. Then we could use the principle that primitive $p$ divides $m$ times $g$, $m$ not equal to $0$, primitive --- we've used this principle (Recall: \cref{prop:primitive-divisibility}) \inlinemetanote{points at the upper board} --- implies $p$ divides $g$, and so $g$ is in $(p)$, and we're done \inlinemetanote{writes \qt{$p \text{ (prim)} \mid m g \text{ ($m \neq 0$)}$}, \qt{$\implies p \mid g$}, and \qt{$g \in (p)$ \checkmark}}.

All we need to show is: we must show that $p$ irreducible --- and now we use the assumption the degree of $p$ is positive --- implies $p$ is primitive \inlinemetanote{writes \qt{We must show} and \qt{$p \text{ irr., } \deg(p) > 0 \implies p \text{ is primitive}$}}. Okay? That's the last step, the last step in this intricate argument: $p$ irreducible, degree positive, it's primitive.

So, we factor $p$ as the content of $p$ times $p^*$ \inlinemetanote{writes: $p = C(p) \, p^*$}, all right? That's the content factorization of $p$. The fact that $p$ is irreducible: one of these two must be units --- $p$ irreducible implies either $C(p)$ or $p^*$ is a unit \inlinemetanote{writes: $p \text{ irr.} \implies C(p) \text{ or } p^* \text{ unit}$}. That's what it means to be irreducible: one of these two is a unit.

This thing \inlinemetanote{points at $p^*$} cannot be a unit --- why is that? Because it has the same degree as $p$, which is positive, and positive-degree things are not units. So, this thing \inlinemetanote{points at $C(p)$} has to be a unit. So, that can't be a unit because it's positive degree, so that has to be a unit, and that means that $p$ is primitive \inlinemetanote{crosses out \qt{$p^*$}, circles \qt{$C(p)$}, writes \qt{$\deg > 0$} under $p^*$ and \qt{unit} under $C(p)$, and writes \qt{$\implies p \text{ is primitive}$}}.

So, at the very last stage, we used the hypothesis that the degree is greater than $0$. And that's the end of the proof.
\end{speech}

\begin{content}[{Primitivity of $p$ and Final Deduction $g \in (p)$}]
To finish from \eqref{eq:p-div-mg}, it suffices to show that $p$ is primitive; then \cref{prop:primitive-divisibility} applied to $p \mid mg$ with $m \in R \setminus \{0\}$ yields $p \mid g$, i.e., $g \in (p)$.

\begin{claim}[Irreducible Polynomials of Positive Degree Are Primitive]
\label[claim]{clm:irr-pos-deg-primitive}
If $p \in R[x]$ is irreducible and $\deg(p) > 0$, then $p$ is primitive.
\end{claim}

\begin{proof}[Proof of \cref{clm:irr-pos-deg-primitive}]
Factor $p$ via its content as
\[
p = C(p) \, p^*,
\]
where $C(p) \in R$ and $p^* \in R[x]$ is primitive. Since $p$ is irreducible in $R[x]$, either $C(p)$ or $p^*$ must be a unit in $R[x]$. However, $\deg(p^*) = \deg(p) > 0$, and no polynomial of positive degree can be a unit in $R[x]$ (since $R$ is an integral domain). Therefore, $p^*$ is not a unit, which forces $C(p)$ to be a unit in $R$. By definition ($C(p) \stackrel{u}{=} 1$), $p$ is primitive.
\end{proof}

Applying \cref{prop:primitive-divisibility} to $p \mid mg$ with $p$ primitive and $m \in R \setminus \{0\}$ now gives $p \mid g$, hence $g \in (p)$.
\end{content}

\begin{content}[Complete Proof Summary: Case (ii) of Gauss's Theorem]
We assemble all steps of Case (ii) into a single continuous argument:

\textbf{1. Setup and Minimal-Degree Element:} Assume $p \in R[x]$ is irreducible with $\deg(p) > 0$, $fg \in (p)$, and $f \notin (p)$. Set $I := (f, p) \subset R[x]$ as in \eqref{eq:ideal-fp}. Since $p \neq 0$, $I \neq (0)$, so we may choose $m \in I \setminus \{0\}$ of minimal degree as in \eqref{eq:minimal-m}.

\textbf{2. Passing to $Q(R)[x]$ and Clearing Denominators:} Viewing $m, p, f, g$ in $Q[x]$ over the fraction field $Q := Q(R)$, the division algorithm in $Q[x]$ gives $f = \hat{q} m + \hat{r}$ \eqref{eq:div-f-qhat} with $\hat{r} = 0$ or $\deg(\hat{r}) < \deg(m)$. Multiplying by a common non-zero denominator $b' \in R \setminus \{0\}$ yields $b' f = q' m + r'$ in $R[x]$ \eqref{eq:cleared-f}, with $r' = 0$ or $\deg(r') < \deg(m)$. Since $r' = b'f - q'm \in I$, minimality of $\deg(m)$ forces $r' = 0$. Factoring $m = C(m) m^*$ gives $b' f = q' C(m) m^*$ \eqref{eq:bprime-f}; since $m^*$ is primitive and $b' \neq 0$, \cref{prop:primitive-divisibility} yields $m^* \mid f$ \eqref{eq:mstar-div-f}.

\textbf{3. Repeating for $p$ and Using Irreducibility:} Applying the exact same division and denominator-clearing step to $p$ gives $b'' p = q'' m = q'' C(m) m^*$ \eqref{eq:bdoubleprime-p}, whence $m^* \mid p$ \eqref{eq:mstar-div-p} and $p = m^* B$ \eqref{eq:p-mstar-B} for some $B \in R[x]$. By irreducibility of $p$, either $m^*$ or $B$ is a unit in $R[x]$. If $B$ were a unit, then $p \mid m^*$ and thus $p \mid f$, contradicting $f \notin (p)$. Hence, $m^*$ is a unit in $R[x]$, so $m^* \in R$, and therefore, $m = C(m) m^* \in R \setminus \{0\}$ \eqref{eq:m-in-R}.

\textbf{4. Concluding $g \in (p)$:} Writing $m = Af + Bp$ for $A, B \in R[x]$ \eqref{eq:m-linear-combo} and multiplying by $g$ yields $mg = Afg + Bgp$ \eqref{eq:mg-combo}. Since $fg \in (p)$, we have $p \mid mg$ \eqref{eq:p-div-mg} with $m \in R \setminus \{0\}$. Finally, factoring $p = C(p) p^*$ and using $\deg(p^*) = \deg(p) > 0$ shows $p^*$ is not a unit, so irreducibility of $p$ forces $C(p)$ to be a unit; thus, $p$ is primitive (\cref{clm:irr-pos-deg-primitive}). By \cref{prop:primitive-divisibility}, $p \mid g$, so $g \in (p)$.

Combining Case (i) ($\deg(p) = 0$) and Case (ii) ($\deg(p) > 0$) shows that every irreducible in $R[x]$ generates a prime ideal; together with existence of factorizations, this completes the proof that $R[x]$ is a UFD (\cref{thm:ufd-polynomial-ring-w7tue}).
\end{content}

\end{proofWrapper}

\section{Reflections on Studying Rings and the Nature of Learning}

\begin{speech}[00:39:17 - 00:40:59]
So, that's the whole argument, and as I said, it has a certain intricacy; it uses most of the things we've discussed about rings. And that, unfortunately, is the last result about rings in this class. That's not really true --- there's a little bit at the end ---, but insofar as the development of the theory of rings, we are finished, so you can bid a tearful goodbye to rings now.

The next topic is groups, but I want to say a few things, because people asked me some questions, and they're good questions. Like, the first question is: how are you supposed to study this material? That's a really good question. \inlinemetanote{begins erasing the lower board}

Another question is: how did I learn the material? That's not such a good question; it's not so relevant, but... Maybe a deeper question, which goes to the heart of the matter, is, you know: what are we doing here? That's a really good question. And if you're studying, you should take some time to think about these questions. There are different ways. So, maybe I try to answer these questions in different --- I have a few minutes, maybe it's worth giving a few.

So, the first question about how to study the material: you can say, what have I taught in this class about rings? And you can say there's kind of three things. One is definitions --- and as I said many, many times, those are incredibly important; it's what makes all these proofs work perfectly. And they're not tautological: at the end we prove some actual non-trivial results, but the definitions are --- there's a lot of mathematical insight and knowledge in the definitions \inlinemetanote{walks to the far left board}.
\end{speech}

\begin{speech}[00:40:59 - 00:41:36]
So, I would say study the definitions, then there's theorems, and then there's examples \inlinemetanote{writes \qt{Def}, \qt{Thms}, and \qt{Examples}}. These are like the three things you should really kind of try to understand \inlinemetanote{circles \qt{Def}, \qt{Examples}, and \qt{Thms} and connects them with lines into a triangle}. I don't know how to say more than that, but they are somehow equal weights in the class.

If you understand the definitions, the theorems --- you know, these definitions and the objects are like chess pieces, and each of them has certain moves, and the theorems involve playing out those moves. So, if you understand this, you'll understand that also. That's like at the lowest level.
\end{speech}

\begin{content}[The Three Pillars of Studying the Material]
\definecolor{figBlue}{HTML}{07608F} % RGB (7, 96, 143): petrol blue (MidnightBlue)
\definecolor{figGray}{HTML}{666665} % RGB (102, 102, 101): dark gray (gray!80!black)
\begin{hidden}
- Canvas: no scale, cm; triangle of three equal-weight pillars drawn on the board: Def, Examples, Thms. Picture spans
  x in [-0.3, 5.9], y in [-0.3, 4.7]
- Order: (1) circle fills and outlines for Def, Examples, Thms; (2) connecting edges between node borders; (3) text last
- Nodes (text at \small, 0.18 per character):
  def: centre (2.8, 3.8), circle minimum size 1.7 cm (radius 0.85); "Def" (3 chars, 0.54 -> 32 % of 1.7)
  ex:  centre (0.8, 0.8), circle minimum size 2.2 cm (radius 1.10); "Examples" (8 chars, 1.44 -> 65 % of 2.2)
  thm: centre (4.8, 0.8), circle minimum size 2.2 cm (radius 1.10); "Thms" (4 chars, 0.72 -> 33 % of 2.2)
- Gaps: ex-thm centre distance 4.0, border gap 4.0 - 2.2 = 1.8 cm; def-ex and def-thm centre distances
  sqrt(2^2 + 3^2) = 3.61, border gaps 3.61 - 0.85 - 1.10 = 1.66 cm.
- Collisions: no labels outside the three nodes; lines terminate cleanly at node boundaries; text drawn last inside nodes.
\end{hidden}
\begin{center}
\begin{tikzpicture}[font=\small]
    % (1) Node circles
    \node[circle, draw=figBlue, thick, fill=figBlue!10, minimum size=1.7cm] (def) at (2.8, 3.8) {};
    \node[circle, draw=figBlue, thick, fill=figBlue!10, minimum size=2.2cm] (ex)  at (0.8, 0.8) {};
    \node[circle, draw=figBlue, thick, fill=figBlue!10, minimum size=2.2cm] (thm) at (4.8, 0.8) {};

    % (2) Connecting edges between the three pillars
    \draw[thick, figGray] (def) -- (ex) (ex) -- (thm) (thm) -- (def);

    % (3) Text last
    \node[figBlue] at (def) {Def};
    \node[figBlue] at (ex)  {Examples};
    \node[figBlue] at (thm) {Thms};
\end{tikzpicture}
\end{center}
\end{content}

\begin{insight}[Definitions as Chess Pieces]
The lecturer captures the relation between the three pillars with a chess analogy: \qt{these definitions and the objects are like chess pieces, and each of them has certain moves, and the theorems involve playing out those moves.} Rather than being mere tautologies, the definitions encode deep mathematical insight, so mastering them unlocks how the theorems operate.
\end{insight}

\begin{speech}[00:41:36 - 00:43:02]
But this question about what are we doing here --- that's an interesting question. And many people have written about it; it's been, you know, $2000$ years of intellectual history about that, and you can go look at it.

One thing that I recommend my students read is an essay by Schopenhauer, which is called \qt{On Thinking for Yourself}, which I like a lot. He's a very dense, repetitive author, so you can read just half the essay; it's a short essay. The German is \qt{Selbstdenken}; it's a nice essay. Most of the students don't like it after they read it, but anyway, it's a really good essay. It's kind of about this topic.

And, you know, all the greats have written about it. If you read the past --- I mean, there's a lot of things to say, I could give a whole lecture ---, but if you read Plato, his \qt{Republic}, he talks about what education is about; Socrates, he didn't write, but there's a lot of stuff there.

The one sentence which I like the best is Confucius --- and you think, is he going to quote Confucius in this class? ---, but I will quote it: I think the statement that best captures what we do is Confucius wrote about teaching, he says --- roughly speaking --- \qt{the truth has $4$ corners; as a teacher I can show you $1$, the other $3$ are for the student to discover themselves.} That's what we're doing here. Okay, that's it. \inlinemetanote{video cuts slightly after a pause for student questions}
\end{speech}

\begin{insight}[{Confucius on Teaching and Independent Discovery}]
To answer \qt{what are we doing here?}, the lecturer points students to Schopenhauer's essay \emph{On Thinking for Yourself} (\emph{Selbstdenken}), Plato's \emph{Republic}, and Confucius's remark on teaching (\emph{Analects}~VII.8): \qt{the truth has $4$ corners; as a teacher I can show you $1$, the other $3$ are for the student to discover themselves.}
\end{insight}

\begin{hidden}
% timestamp: 00:43:02
% topic: End of ring theory, study advice (Def/Examples/Theorems), and philosophical reflection on learning
% board_state: triangle diagram (Def, Examples, Thms)
% next_goal: Address student question on where deg(p) > 0 was used in Case (ii), then begin Groups
% open_loops: none
\end{hidden}

\subsection{Addendum: Where \texorpdfstring{$\deg(p) > 0$}{deg(p) > 0} Was Used}

\begin{speech}[00:43:02 - 00:44:10]
Okay, so, there were some questions about the proof. The question, roughly speaking, was: in that whole hour, I was assuming the degree of $p$ is positive \inlinemetanote{writes: $\deg(p) > 0$}: where do I use it?

And it was just --- it's like a player who sits on the bench for the whole game, but at the end we use this: $p$ irreducible, degree of $p$ positive implies $p$ primitive \inlinemetanote{writes: $p \text{ irr., } \deg(p) > 0 \implies p \text{ is primitive}$}. That was what we needed at the last moment. So, this player just sits on the bench, at the last moment scores the winning goal. That's the only time it's used, at the end \inlinemetanote{circles the implication}. You can look at the argument, but it's the only time it's used.

Or... yeah, so if anyone reads Schopenhauer, you should tell me.

Oh yeah, and finally the question about how did I learn it: the answer is, at a certain point in my life I studied very hard. That's how you learn it. That's the only way. So, it's not very helpful, I think. Okay, so...
\end{speech}

\begin{content}[{Exact Point Where $\deg(p) > 0$ Enters the Proof}]
Throughout Case (ii) of \cref{thm:ufd-polynomial-ring-w7tue}, the hypothesis $\deg(p) > 0$ is used in only one place, at the very end (\cref{clm:irr-pos-deg-primitive}):
\[
p \text{ irreducible}, \; \deg(p) > 0 \implies p \text{ is primitive}.
\]
Specifically, when factoring $p = C(p) p^*$, $\deg(p^*) = \deg(p) > 0$ ensures that $p^*$ cannot be a unit in $R[x]$, forcing $C(p) \in \mathcal{U}_R$.
\end{content}

\begin{insight}[{The \qt{Player on the Bench} Analogy}]
The lecturer compares the hypothesis $\deg(p) > 0$ to \qt{a player who sits on the bench for the whole game} and \qt{at the last moment scores the winning goal}: every step up to $p \mid mg$ holds even if $\deg(p) = 0$, and $\deg(p) > 0$ is needed solely to deduce that $p$ is primitive from its irreducibility.
\end{insight}

\section{Introduction to Groups}

\subsection{Group Axioms and Elementary Properties}

\begin{speech}[00:44:10 - 00:45:56]
Yeah, so, now we switch to groups \inlinemetanote{writes and underlines \qt{Groups}}. And we've already seen the definition, so I don't really say so much about it (Recall: \cref{def:group}). A group is... first: Groups --- that's the topic now \inlinemetanote{underlines \qt{Groups} again} ---, so $G$ with a binary operation and a neutral element is a group \inlinemetanote{writes: $(G, \cdot, e) \text{ is a group}$}.

Here we have associativity, the unit property, and the existence of inverse \inlinemetanote{writes \qt{$\bullet \text{ Associativity}$}, \qt{$\bullet \text{ Unit property}$}, and \qt{$\bullet \text{ Existence of inv.}$}}. So, we've seen all of this; I'm not going to say so much about it except mostly about notation. Associativity --- they're all interesting and they're worth thinking about again now, because I defined all of this maybe at the beginning of the semester, and you've probably seen the definition of a group even before.

So, associativity is something that says that if I do the binary operation in two different ways, it's the same \inlinemetanote{writes: $(x \cdot y) \cdot z = x \cdot (y \cdot z)$}. We haven't said anything about associativity also for rings, so maybe I say a moment --- I'll say in a moment.

The unit property says that there exists a unit such that $ex = xe = x$ for all $x$ in the group \inlinemetanote{writes: $\exists\, e \text{ s.t. } ex = xe = x \;\; \forall x \in G$}.

And the existence of inverse says that for all $x$ in the group, there exists $y$ such that $xy = yx$ equals the unit \inlinemetanote{writes: $\forall x \in G, \; \exists\, y \text{ s.t. } xy = yx = e$}.
\end{speech}

\begin{content}
\recap{Groups were defined in Week~2 (\cref{def:group}), together with uniqueness of the neutral element and of inverses (\cref{prop:uniqueness-neutral-element}, \cref{prop:uniqueness-inverse}). The lecturer reviews them below.}

\begin{definition}[Group]\label[definition]{def:group-w7tue}
A tuple $(G, \cdot, e)$, consisting of a set $G$, a binary operation $\cdot$, and a neutral element $e \in G$, is a \newterm{group} if it satisfies the following three axioms:
\begin{itemize}
    \item \textbf{(G1) Associativity:} $(x \cdot y) \cdot z = x \cdot (y \cdot z)$ for all $x, y, z \in G$.
    \item \textbf{(G2) Unit property:} $\exists\, e \in G$ such that $ex = xe = x \quad \forall x \in G$.
    \item \textbf{(G3) Existence of inverses:} $\forall x \in G, \; \exists\, y \in G$ such that $xy = yx = e$.
\end{itemize}
\end{definition}
\end{content}

\begin{speech}[00:45:56 - 00:46:33]
And there are many things to say, and I leave this to you --- you can say that, you know, the unit is unique \inlinemetanote{writes: $\leftarrow \text{Unit is unique}$}. I will say one thing here, because unit... \inlinemetanote{puts parentheses around \qt{Unit}, giving \qt{(Unit) is unique}}

These proofs are all very simple. So, the idea is that if we have both --- these are both units \inlinemetanote{opens a parenthesis and writes: $e, \hat{e} \text{ units in } G$} ---, this is a simple proof, then I just take $e$ times $\hat{e}$ \inlinemetanote{writes: $e \hat{e}$}. Using the unit property of $e$, I get this is equal to $\hat{e}$, and using the unit property of $\hat{e}$, I get this is equal to $e$ \inlinemetanote{writes \qt{$= \hat{e}$}, writes a vertical equals sign below $e\hat{e}$ with $e$ beneath it, and then closes the parenthesis}. So, this shows the unit is unique.
\end{speech}

\begin{content}
\begin{proposition}[Uniqueness of the Unit]\label[proposition]{prop:group-unit-unique}
In a group $(G, \cdot, e)$, the unit (neutral element) is unique.
\end{proposition}

\begin{proof}
Suppose $e$ and $\hat{e}$ are both units in $G$, and consider the product $e \hat{e}$. Since $e$ is a unit in $G$, we have $e \hat{e} = \hat{e}$; since $\hat{e}$ is a unit in $G$, we have $e \hat{e} = e$. Combining the two equalities gives $e = e \hat{e} = \hat{e}$. Thus, the unit is unique.
\end{proof}
\end{content}

\begin{speech}[00:46:33 - 00:47:40]
Something that we use a lot in groups is left and right cancellation \inlinemetanote{writes and underlines: $\text{Left / right cancellation}$}. This says that if I have an equation $xa = xb$, then I can cancel $x$: this implies $a$ is equal to $b$ \inlinemetanote{writes \qt{$\text{if } xa = xb$} and \qt{$\implies a = b$}}. In a group you can always cancel.

And the reason is because I take this equation $xa = xb$, and I multiply by this inverse of $x$ \inlinemetanote{writes: $xa = xb$}. I can then use $y$ is equal to the inverse of $x$ \inlinemetanote{writes: $\text{use } y = \text{inv. of } x$}, and I act on this: I get $yxa = yxb$ \inlinemetanote{writes: $yxa = yxb$}. And then these are units: it's the neutral element, and this implies $a = b$ \inlinemetanote{writes parentheses around $(yx)$ on both sides, writes \qt{$ea = eb$} below, and writes \qt{$\implies a = b$}}.

So, we'll use this all the time, also on the other side. That's the left and right cancellation property.
\end{speech}

\begin{content}
\begin{proposition}[Left and Right Cancellation]\label[proposition]{prop:cancellation-laws}
Let $(G, \cdot, e)$ be a group and $x, a, b \in G$.
\begin{enumerate}[label=\textbf{(\arabic*)}]
    \item \textbf{Left cancellation:} If $xa = xb$, then $a = b$.
    \item \textbf{Right cancellation:} If $ax = bx$, then $a = b$.
\end{enumerate}
\end{proposition}

\begin{proof}
\textbf{Left cancellation:} Suppose $xa = xb$. Let $y \in G$ be an inverse of $x$ (\textbf{(G3)}), so that $yx = e$. Multiplying both sides on the left by $y$ gives:
\begin{align*}
yxa &= yxb \\
(yx)a &= (yx)b && \text{(by associativity \textbf{(G1)})} \\
ea &= eb && \text{(since $yx = e$)} \\
a &= b && \text{(by the unit property \textbf{(G2)}).}
\end{align*}

\textbf{Right cancellation:} Analogously, if $ax = bx$, multiplying both sides on the right by an inverse $y \in G$ of $x$ (using $xy = e$) gives $(ax)y = (bx)y$, and hence, $a(xy) = b(xy)$ by \textbf{(G1)}, $ae = be$, and $a = b$ by \textbf{(G2)}.
\end{proof}
\end{content}

\begin{speech}[00:47:40 - 00:48:15]
The final thing is --- I'll say it here, I've said all these things before, but maybe just some... the inverse is unique \inlinemetanote{draws a large opening parenthesis above the unit-uniqueness proof, writes and underlines: $\text{Inv of } x \text{ is unique}$}.

That means if I have $xy$ is equal to $e$, and this also equals $x\hat{y}$ \inlinemetanote{writes: $xy = e = x\hat{y}$} --- so, that's like $y$ and $\hat{y}$ are now two putative inverses ---, then I get $xy$ is equal to $x\hat{y}$, and then I use the cancellation property (Recall: \cref{prop:cancellation-laws}) \inlinemetanote{writes \qt{$xy = x\hat{y}$} and \qt{$\implies y = \hat{y}$}, closes the parenthesis, and points at the cancellation rule}.

Okay, so, those are the really basic aspects of it.
\end{speech}

\begin{content}
\begin{proposition}[Uniqueness of Inverses]\label[proposition]{prop:group-inverse-unique}
For each element $x$ in a group $(G, \cdot, e)$, its inverse is unique.
\end{proposition}

\begin{proof}
Suppose $y, \hat{y} \in G$ are both inverses of $x$. In particular,
\[
xy = e = x\hat{y}.
\]
Hence, $xy = x\hat{y}$, and left cancellation (\cref{prop:cancellation-laws}) immediately implies $y = \hat{y}$.
\end{proof}
\end{content}

\begin{aiNote}[Beyond the Lecture: A Second Proof of Uniqueness of Inverses]
We also saw a direct associativity proof of uniqueness of inverses in Week~2 (\cref{prop:uniqueness-inverse}) that avoids appealing to cancellation separately: if $y$ and $\hat{y}$ are both inverses of $x$, then $y = ye = y(x\hat{y}) = (yx)\hat{y} = e\hat{y} = \hat{y}$. The cancellation argument given here packages the exact same computation into \cref{prop:cancellation-laws}.
\end{aiNote}

\begin{hidden}
% timestamp: 00:48:15
% topic: Group definition and basic properties (uniqueness of unit, left/right cancellation, uniqueness of inverses)
% board_state: def:group-w7tue, prop:group-unit-unique, prop:cancellation-laws, prop:group-inverse-unique
% next_goal: Discuss focus on finite non-commutative groups, ring additive groups, notation, and preview S_n
% open_loops: associativity for n-fold products (mentioned at 00:45:24)
\end{hidden}

\subsection{Focus of the Course, Notation, and First Examples}

\begin{speech}[00:48:15 - 00:50:18]
So, about groups: in some sense it's algebra just like about rings --- rings have more structure. We were very much interested, because of the classical examples, primarily in rings that have infinite numbers of elements --- not exclusively, but primarily, our classical examples. Our focus on groups will be primarily on groups with finite numbers of elements \inlinemetanote{pushes the upper board down, moves to the new lower board, and writes: $\longrightarrow \text{Groups with finite numbers of elements}$}.

Groups can have, of course, infinite numbers of elements. So, we will primarily be interested in groups with finite numbers of elements. That's a different perspective than our study of rings, where I would say we were primarily interested in rings which were infinite rings.

And we will primarily be interested in groups that are not commutative \inlinemetanote{writes \qt{$\longrightarrow \text{groups which are not commutative}$} and underlines \qt{not} twice}. So, that's another difference.

We have now for this period about groups to be extremely careful never to switch order. I was cavalier about that; I never cared about the order because our rings were commutative, but now when we apply the group operation, you have to be very careful about the order. You're not allowed to switch the order unless you have a special reason for it.

Okay. And the examples, of course, the examples we've seen before: if you take $R$ is any ring --- that's a ring ---, then you can make a group \inlinemetanote{writes \qt{$(R, +, \cdot, 0, 1)$} and \qt{$\longrightarrow (R, +, 0) \text{ Grp}$}}. And that was our definition of a ring (Recall: \cref{def:ring}): the additive structure is always a group. These are abelian groups; they're commutative groups. They won't be primarily what we're interested in, but that's a solid example of a group.
\end{speech}

\begin{content}[Scope of Group Theory in the Course and Additive Ring Groups]
Unlike the preceding study of rings --- which focused primarily on infinite, commutative rings (\cref{rem:ring-convention}) --- the study of groups in this course will focus primarily on:
\begin{itemize}
    \item groups with finite numbers of elements, and
    \item groups which are \emph{not} commutative (so the order of factors in a product may never be switched without justification).
\end{itemize}

\begin{example}[Additive Group of a Ring]\label[example]{ex:ring-additive-group}
Every ring $(R, +, \cdot, 0, 1)$ (\cref{def:ring}) gives rise to an abelian (commutative) group under addition by forgetting multiplication:
\[
(R, +, \cdot, 0, 1) \longrightarrow (R, +, 0) \quad \text{(a group)}.
\]
\end{example}
\end{content}

\begin{speech}[00:50:18 - 00:51:50]
And you see here the notation for the group; it's a good point here. I said $G$ and dot is the standard for the operation in the group, and $e$ is standard for the neutral element \inlinemetanote{writes: $(G, \cdot, e)$}. The neutral element is also sometimes denoted by $1$, because the dot kind of suggests multiplication, and $1$, psychologically, is helpful --- I don't know \inlinemetanote{draws a downward arrow from $e$ and writes \qt{$1$}}. Sometimes people will write $1_G$ \inlinemetanote{writes \qt{$1_G$} below $1$}. These are just accepted terms when you use dot.

But you can name this anything you want; you could also name it plus. If you name it plus, usually the neutral element is called $0$ \inlinemetanote{writes: $(G, +, 0)$}. But this is just a question about names: you can name this binary operation anything you want and the neutral element anything you want. It's just, you know, if you do too crazy choices you'll confuse everybody. But these are kind of the standard --- so, that's why when I write plus, $0$ is the neutral element.

So, there's lots of examples here. The example that's more relevant for the way we think about groups --- the interesting things we're going to prove about groups --- is the \newterm{symmetric group}, $S_n$; that's the permutation group \inlinemetanote{writes: $S_n = \text{permutation group on } n \text{ elements}$}. I'm going to review this. I think I mentioned this in passing. And in some sense this is like the most interesting finite group --- I will make that claim.

That's... and you should think about these --- you know, I already presented the definitions earlier, but you should think about them.
\end{speech}

\begin{content}
\begin{notation}[Multiplicative and Additive Group Notation]\label[notation]{not:group-notation}
\begin{itemize}
    \item In general multiplicative notation, a group is written $(G, \cdot, e)$; the neutral element $e$ is also commonly denoted $1$ or $1_G$.
    \item When the group operation is written additively, the group is denoted $(G, +, 0)$, with neutral element $0$.
\end{itemize}
\end{notation}

A central example of a finite group is the \emph{symmetric group}:
\[
S_n = \text{permutation group on } n \text{ elements},
\]
which will be defined and studied formally below (\cref{def:symmetric-group}).
\end{content}

\subsection{Generalized Associativity and Catalan Numbers}

\begin{speech}[00:51:50 - 00:53:41]
And there's a lot of interesting things here which we don't think --- like I didn't talk at all about this associativity property. And in some sense in your education it's kind of late to talk about associativity, because you've been, you know, using operations which are associative since childhood, and not thinking about them.

So, maybe take a moment for, like, mathematical culture: if I have this product of $4$ elements in the group \inlinemetanote{writes: $x_1 \cdot x_2 \cdot x_3 \cdot x_4$}, we just write it like this, like if you're multiplying numbers, etc. But what do you mean by this? Because you're only allowed to do $2$ elements at a time.

One thing you could mean is this: you do these two, get an answer, and do these two, get an answer, and that is equal to your answer \inlinemetanote{writes: $(x_1 \cdot x_2) \cdot (x_3 \cdot x_4) = ?$}. That makes sense, because I do these two at a time and those two at a time, right?

Or I could do these two, and then I could do that one, and then I could send it to that one at the end \inlinemetanote{writes: $((x_1 \cdot x_2) \cdot x_3) \cdot x_4 = ?$}. That's a different way of doing it.

And the associativity axiom says that these two are the same, which you prove by induction --- which probably you've never done, or I don't know whether you've done it. Unless you're really deeply paranoid, you probably have never done that. But it might be fun to do it.

And if you do it, you discover some beautiful things. So --- this is now nothing to do with the class, but mathematical culture ---, you can ask: how many different ways can I put these brackets \inlinemetanote{writes: $\text{How many ways can}$ / $\text{put these brackets}$}? Meaning if I have $x_1 \cdots x_n$ \inlinemetanote{writes: $x_1 \cdots x_n$}, I can group them --- I can do this operation in many different ways. How many ways?
\end{speech}

\begin{content}[Parenthesizing Multi-Factor Products]
Because a group operation is binary, an unparenthesized $4$-fold product
\[
x_1 \cdot x_2 \cdot x_3 \cdot x_4
\]
can be evaluated via different binary groupings, such as
\[
(x_1 \cdot x_2) \cdot (x_3 \cdot x_4) \qquad \text{or} \qquad \bigl((x_1 \cdot x_2) \cdot x_3\bigr) \cdot x_4.
\]
By induction using the associativity axiom, all valid parenthesizations of a product yield the same group element.

\begin{question}[Number of Parenthesizations]\label[question]{q:catalan-brackets}
How many ways can brackets be placed to evaluate an $n$-fold product $x_1 \cdots x_n$ using a binary operation?
\end{question}
\end{content}

\begin{speech}[00:53:41 - 00:55:19]
And that has an extremely famous answer, which is part of math culture, which I am going to write. The hint to this is --- I'm kind of really bad at languages, but anyway, it's this \inlinemetanote{writes: $\text{u, u, dos, cinc, catorze, \dots}$}: maybe like that. Does anyone recognize these numbers? No one?

Aren't there any Catalan speakers in this room? Does anyone speak Catalan? You speak Catalan? \inlinemetanote{student responds inaudibly} So, you don't recognize if I got it right. \inlinemetanote{student responds inaudibly} Spanish --- Spanish is close enough, but...

You seriously don't recognize these numbers? Well, this was supposed to be a huge hint: $1, 1, 2, 5, 14$ \inlinemetanote{writes $1, 1, 2, 5, 14$ underneath \qt{u, u, dos, cinc, catorze}}. So, these are the Catalan numbers! That's actually Catalan --- that's my attempt to write them in Catalan. They don't come from the language; there's a French mathematician named Catalan.

The formula for this is $C_n$ is equal to $\frac{1}{n+1} \binom{2n}{n}$ \inlinemetanote{writes: $C_n = \frac{1}{n+1}\binom{2n}{n}$}. And these are some of the most famous numbers in mathematics --- which doesn't mean... they're some of the most famous numbers in mathematics, but they probably don't occur in any course, I don't know. So, that's the $n$-th Catalan number. It's a totally beautiful object.

And if you solve this problem correctly, the number of ways to bracket $n$ things, this is equal to $C_{n-1}$ \inlinemetanote{writes \qt{$C_{n-1}$} below $x_1 \cdots x_n$}: it's the $(n-1)$-st Catalan number.

Anyway. So, that's already just trying to unpack associativity, what it means. If you do it, and you follow your nose, eventually you get to the Catalan numbers.
\end{speech}

\begin{content}[{Catalan Numbers and Answer to \cref{q:catalan-brackets}}]
\begin{answer}[to \cref{q:catalan-brackets}]\label[answer]{ans:catalan-brackets}
The number of ways to parenthesize an $n$-fold product $x_1 \cdots x_n$ is the $(n-1)$-st \newterm{Catalan number} $C_{n-1}$, where the sequence of Catalan numbers starts with
\[
1, \; 1, \; 2, \; 5, \; 14, \dots
\]
(written on the board in Catalan as \qt{u, u, dos, cinc, catorze, \dots}) and has the closed-form formula
\begin{equation}\label{eq:catalan-formula}
C_n = \frac{1}{n+1} \binom{2n}{n} \qquad (n \ge 0).
\end{equation}
\end{answer}

\begin{steps}[Indexing Check for Small Numbers of Factors]
With index $k = 0, 1, 2, 3, 4$, \eqref{eq:catalan-formula} gives $C_0 = 1$, $C_1 = 1$, $C_2 = \frac{1}{3}\binom{4}{2} = 2$, $C_3 = \frac{1}{4}\binom{6}{3} = 5$, and $C_4 = \frac{1}{5}\binom{8}{4} = 14$.
\begin{itemize}
    \item For $3$ factors $x_1 x_2 x_3$, there are $C_2 = 2$ bracketings: $(x_1 x_2)x_3$ and $x_1(x_2 x_3)$.
    \item For $4$ factors $x_1 x_2 x_3 x_4$, there are $C_3 = 5$ bracketings: the two shown earlier, plus $x_1\bigl((x_2 x_3) x_4\bigr)$, $x_1\bigl(x_2 (x_3 x_4)\bigr)$, and $\bigl(x_1(x_2 x_3)\bigr) x_4$.
\end{itemize}
\end{steps}
\end{content}

\begin{aiNote}[Beyond the Lecture: Generalized Associativity]
% [PERSISTENT, ADDITION]
Textbooks call the statement that every valid parenthesization of $x_1 \cdots x_n$ yields the same product the \emph{generalized associative law}.
\end{aiNote}

\begin{aiNote}[Beyond the Lecture: Catalan Numbers and Eugène Charles Catalan]
% [PERSISTENT, ADDITION]
The Catalan numbers are named after the Belgian/French mathematician Eugène Charles Catalan (1814--1894). In combinatorics, $C_n$ also counts the number of Dyck paths of length $2n$, rooted binary trees with $n$ internal vertices, and triangulations of a convex polygon with $n+2$ vertices.
\end{aiNote}

\begin{speech}[00:55:19 - 00:56:29]
But have you seen the definition of the Catalan numbers before? Yeah, that's a tragedy.

\inlinemetanote{erases the lower chalkboard} But I don't think they occur in any class, I don't know; maybe they do, hard to say. You can't actually learn mathematics by going to class; that's the other problem.

So, the number of ways you can do $3$ things is $2$, and the number of ways you can do $4$ things is $5$. That's what they are in the...
\end{speech}

\begin{hidden}
% timestamp: 00:56:29
% topic: Course focus on finite non-commutative groups, notation, and Catalan numbers
% board_state: ex:ring-additive-group, not:group-notation, q:catalan-brackets, ans:catalan-brackets, eq:catalan-formula
% next_goal: Frame group inverse as a map, define subgroups and group homomorphisms, prove f(e_1) = e_2
% open_loops: none
\end{hidden}

\subsection{Inverse Map, Subgroups, and Group Homomorphisms}

\begin{speech}[00:56:29 - 00:57:52]
Okay, so, then there's other things that also follow from what I already explained. If I have a group \inlinemetanote{writes: $G$} --- another way to think about the operations: there's this binary operation \inlinemetanote{writes: $\bullet \text{ binary}$}, and we have the neutral element \inlinemetanote{writes: $e \in G$}, and it's nice to think also about inverse being a map: inverse from the group to the group \inlinemetanote{writes: $\operatorname{Inv}\colon G \to G$}. Meaning if I have $x$, I have a unique inverse $x^{-1}$ \inlinemetanote{writes: $x \mapsto x^{-1}$}. So, it's a function from the group to the group; it takes an element to its inverse.

All right. And so there's a notion of a \newterm{subgroup} \inlinemetanote{writes and underlines \qt{Subgroup}, then writes: $S \subset G$}. So, it's a subset, and it should be closed under the operation, meaning if I take two elements of the subset and I use the group operation I stay in it; it should contain the neutral element; and one more thing is $S$ should have inverses, the inverse maps $S$ to $S$ \inlinemetanote{writes and circles \qt{$(S, \cdot, e)$}, and writes \qt{$\cdot \colon S \times S \to S$}, \qt{$e \in S$}, and \qt{$\operatorname{Inv} \colon S \to S$}}. So, every element in $S$ has its inverse in $S$. That's what a subgroup is. Okay?
\end{speech}

\begin{content}
\recap{Subgroups and group homomorphisms were defined in Week~2 (\cref{def:subgroup}, \cref{def:group-homomorphism}), where $f(e_1) = e_2$ was also shown (\cref{prop:hom-preserves-identity}).}

\begin{remark}[The Inverse Map on a Group]\label[remark]{rem:inverse-map}
Because inverses in a group $G$ are unique (\cref{prop:group-inverse-unique}), inversion defines a function $\operatorname{Inv} \colon G \to G$ sending $x \mapsto x^{-1}$. Thus, a group $G$ comes equipped with a binary operation, a neutral element $e \in G$, and the map $\operatorname{Inv} \colon G \to G$.
\end{remark}

\begin{definition}[Subgroup]\label[definition]{def:subgroup-w7tue}
A subset $S \subset G$ (throughout this course, $\subset$ allows equality, i.e., $\subseteq$) of a group $(G, \cdot, e)$ is a \newterm{subgroup} (so $(S, \cdot, e)$ is a group under the restricted operation) if:
\begin{enumerate}[label=\textbf{(\arabic*)}]
    \item $S$ is closed under multiplication: $\cdot \colon S \times S \to S$,
    \item $S$ contains the neutral element: $e \in S$,
    \item $S$ is closed under inverses: $\operatorname{Inv} \colon S \to S$ (i.e., $x^{-1} \in S$ for all $x \in S$).
\end{enumerate}
\end{definition}
\end{content}

\begin{speech}[00:57:52 - 00:58:47]
And we have a notion of \newterm{group homomorphisms} \inlinemetanote{writes: $\text{Homs:}$}. If $G_1$ and $G_2$ --- well, let me write them out: $G_1$ with a neutral element and $G_2$ are groups \inlinemetanote{writes: $(G_1, \cdot, e_1), (G_2, \cdot, e_2)$ are groups}, then $f$ from $G_1$ to $G_2$ is --- this is a definition --- a group hom, group homomorphism \inlinemetanote{writes \qt{$f \colon G_1 \to G_2$}, writes \qt{is a group hom. if}, and underlines \qt{group hom.}}. When I underline things, I always mean definition. It is a group homomorphism if $f(xy)$ is equal to $f(x) f(y)$ \inlinemetanote{writes: $f(xy) = f(x) \cdot f(y)$}.

So, here is where I use the group operation on $G_1$ \inlinemetanote{points at $xy$}, and here I use the group operation on $G_2$ \inlinemetanote{points at $f(x) \cdot f(y)$}. So, that's the whole definition --- just this condition.
\end{speech}

\begin{content}
\begin{definition}[Group Homomorphism]\label[definition]{def:group-homomorphism-w7tue}
Let $(G_1, \cdot_1, e_1)$ and $(G_2, \cdot_2, e_2)$ be groups. A map $f \colon G_1 \to G_2$ is a \newterm{group homomorphism} if for all $x, y \in G_1$:
\[
f(xy) = f(x) \cdot f(y),
\]
where the product $xy$ on the left is computed using $\cdot_1$ in $G_1$ and the product $f(x) \cdot f(y)$ on the right is computed using $\cdot_2$ in $G_2$ (explicitly, $f(x \cdot_1 y) = f(x) \cdot_2 f(y)$).
\end{definition}
\end{content}

\begin{speech}[00:58:47 - 01:00:00]
So, you could wonder what happens to the neutral element: so, what happens to $e_1$ \inlinemetanote{writes: $\text{What happens to } e_1?$}?

I just plug it in: $f(e_1 e_1)$ has to be $f(e_1) f(e_1)$ \inlinemetanote{writes: $f(e_1 e_1) = f(e_1) \cdot f(e_1)$}, right? But this thing, by using the unit property, is equal to $f(e_1)$ \inlinemetanote{writes \qt{$f(e_1)$} beneath $f(e_1 e_1)$}, because $e_1 e_1$ is $e_1$, right?

And then I can use cancellation. So, I look at this equation, I just write it: $f(e_1) = f(e_1) \cdot f(e_1)$ \inlinemetanote{writes: $f(e_1) = f(e_1) \cdot f(e_1)$}. So, this is an equation, group equation in $G_2$, right \inlinemetanote{writes: $\in G_2$}? Everything's in $G_2$.

But now I can cancel --- if I want, for this cancellation, I could then put an $e_2$ here \inlinemetanote{inserts \qt{$\cdot e_2$} after $f(e_1)$ on the left}, because that uses the identity in $G_2$, and then I can cancel \inlinemetanote{slashes out $f(e_1)$ on the left and the first $f(e_1)$ on the right, then writes and boxes: $e_2 = f(e_1)$}.

So, that's the answer: what happens to $e_1$? What happens to the neutral element: it has to go to the neutral element in the second one. Okay? It's implied by this condition \inlinemetanote{points at $f(xy) = f(x) \cdot f(y)$}. I don't have to impose it as an additional condition --- I could, it would be harmless to do it, but I don't have to do it.
\end{speech}

\begin{speech}[01:00:00 - 01:01:12]
Then I leave for you to check yourself what happens to $f$ of an inverse \inlinemetanote{writes: $\text{What happens to } f(x^{-1})?$}. And the answer --- it's just the same argument --- is that $f$ of an inverse is equal to the inverse here \inlinemetanote{writes and boxes: $\text{Ans: } f(x^{-1}) = f(x)^{-1}$}. I leave this for you to check, okay? It's the same argument; it respects the inverse operation also.

So, we're going to study the kernels and images just as before for groups, and that'll lead to a theorem which is parallel to what we called the First Isomorphism Theorem (Recall: \cref{thm:first-isomorphism-theorem-statement}) --- but we're not doing that today.

I thought I would do some example, think about the symmetric group, because it's totally crucial to understand that for your future mathematical life. But the symmetric group must have come up in many classes of yours, I would guess. But I won't assume you know anything about it. \inlinemetanote{begins erasing the lower board}
\end{speech}

\begin{content}
\begin{proposition}[Preservation Properties of Homomorphisms]\label[proposition]{prop:hom-preserves-structure}
If $f \colon G_1 \to G_2$ is a group homomorphism, then $f$ automatically preserves the neutral element and inverses:
\begin{enumerate}[label=\textbf{(\arabic*)}]
    \item $f(e_1) = e_2$,
    \item $f(x^{-1}) = f(x)^{-1}$ for every $x \in G_1$.
\end{enumerate}
\end{proposition}

\begin{proof}
Applying the homomorphism property to $e_1 e_1 = e_1$ in $G_1$ gives:
\[
f(e_1) = f(e_1 e_1) = f(e_1) \cdot f(e_1) \quad \text{in } G_2.
\]
Using the unit property of $e_2$ in $G_2$, we rewrite the left-hand side as $f(e_1) \cdot e_2$:
\[
f(e_1) \cdot e_2 = f(e_1) \cdot f(e_1).
\]
By left cancellation in $G_2$ (\cref{prop:cancellation-laws}), we cancel $f(e_1)$ on both sides to obtain
\[
e_2 = f(e_1).
\]
This proves (1). The proof of (2) is left to the students.
\end{proof}

\begin{proof}[Preservation of Inverses (beyond the lecture)]
The lecture leaves part (2) to the students (\qt{I leave for you to check yourself}). Since $x x^{-1} = x^{-1} x = e_1$ in $G_1$, applying $f$ and using part (1) yields:
\[
f(x) \cdot f(x^{-1}) = f(x x^{-1}) = f(e_1) = e_2
\]
and likewise $f(x^{-1}) \cdot f(x) = f(x^{-1} x) = f(e_1) = e_2$. By uniqueness of inverses in $G_2$ (\cref{prop:group-inverse-unique}), $f(x^{-1}) = f(x)^{-1}$.
\end{proof}
\end{content}

\begin{hidden}
% timestamp: 01:01:35
% topic: Inverse map, subgroups, group homomorphisms (preservation of identity and inverses)
% board_state: rem:inverse-map, def:subgroup-w7tue, def:group-homomorphism-w7tue, prop:hom-preserves-structure
% next_goal: Define [n], symmetric group S_n, composition order, arrow diagrams for S_1, S_2, S_3, and |S_n| = n!
% open_loops: kernels, images, and First Isomorphism Theorem for groups (deferred to future lectures)
\end{hidden}

\subsection{The Symmetric Group \texorpdfstring{$S_n$}{S\_n}}

\begin{speech}[01:01:35 - 01:03:34]
So, if we start formally with a discussion of the symmetric group --- and this is, in the trichotomy I said about the subject, this is not definition, this is not theorem, it's example, but it's an incredibly important example: the \newterm{symmetric group} \inlinemetanote{writes and underlines: $\text{Symmetric group:}$}.

The first symbol, which is kind of standard, is square brackets $n$ to be the set of numbers from $1$ to $n$ \inlinemetanote{writes: $[n] = \{1, \dots, n\}$}. It's the $n$-element set, so it's exactly the integers from $1$ to $n$.

And the symmetric group $S_n$ equals the set of functions $f$ from $[n]$ to $[n]$ such that $f$ is bijective \inlinemetanote{writes: $S_n = \{ f \colon [n] \to [n] \mid f \text{ is bijective} \}$}. That's as a set. So, it's just the set of bijections; I hope that's clear.

I need to tell you the operation, and the operation is composition of functions, and I have to give you the neutral element: it's the identity \inlinemetanote{writes and circles $(S_n, \circ, \operatorname{Id})$, then adds a brace and \qt{$\text{def}$} beneath}. That's the identity. There's a unique identity, which takes $1$ to $1$, $2$ to $2$, so on. And that's the definition of the symmetric group. That's the full definition, and you have to think about it, make sure you understand.

One of the important things to think about is what this composition means: if I take $f$ composed with $g$, where $g$ is from $[n]$ to $[n]$ and $f$ is from $[n]$ to $[n]$ \inlinemetanote{writes \qt{$f \circ g$}, with \qt{$g \colon [n] \to [n]$} and \qt{$f \colon [n] \to [n]$} below}. This means --- so, this acts on some number, and the answer is you do $g$ first, and then you do $f$ \inlinemetanote{writes \qt{$(i)$} after $f \circ g$ and writes \qt{$= f(g(i))$}}. That's the definition of composition, right?

Even though, you know, you can confuse yourself by this, because you see $f$ first, nevertheless, the mathematical convention for this symbol \inlinemetanote{points at $f \circ g$} is you do $g$ first, then you do $f$. Okay?
\end{speech}

\begin{content}
\begin{notation}[Standard $n$-Element Set]\label[notation]{not:n-element-set}
For a positive integer $n$, we denote by $[n]$ the standard $n$-element set of integers from $1$ to $n$:
\[
[n] := \{1, \dots, n\}.
\]
\end{notation}

\begin{definition}[{The Symmetric Group $S_n$}]\label[definition]{def:symmetric-group}
As a set, the \newterm{symmetric group} $S_n$ on $n$ elements is the set of all bijections from $[n]$ to itself:
\[
S_n := \{ f \colon [n] \to [n] \mid f \text{ is bijective} \}.
\]
Equipped with composition of functions $\circ$ and the identity map $\operatorname{Id}$, the tuple
\[
(S_n, \circ, \operatorname{Id})
\]
is the symmetric group on $n$ elements.
\end{definition}

\begin{remark}[The Right-to-Left Convention of Composition]\label[remark]{rem:composition-order}
For functions $g \colon [n] \to [n]$ and $f \colon [n] \to [n]$, composition $f \circ g$ is evaluated right-to-left: first $g$, then $f$ (for example, on $1$, and in general $(f \circ g)(k) = f(g(k))$ for $k \in [n]$):
\[
(f \circ g)(1) = f(g(1)).
\]
\end{remark}
\end{content}

\begin{speech}[01:03:34 - 01:04:54]
So, that's the full definition of the group. And you could say: have I checked the axioms? Well, you're supposed to understand that composition of bijective functions gives you a bijective function; you have to also say that the bijective functions have an inverse; and you have to show that composition is associative. I leave all of this to you --- none of those are surprises.

That's the definition of $S_n$, but then there's a question about how to think about it, and that's an interesting question; it's a question about names. How are we going to talk about elements of $S_n$?

So, the primitive way is, say, if you're taking $S_1$, there's nothing to do \inlinemetanote{writes: $S_1 = \{1 \to 1\}$}: $S_1$ has one element, because it's the number of bijections from $1$ to $1$. So, you could say: okay, I could say here's $1$, and here's the map from $1$ to $1$, and it's a bijection \inlinemetanote{draws a diagram mapping $1 \to 1$}. It has one element --- it's not so exciting.

$S_2$, we have bijections from the set of $2$ elements, and one bijection is the identity \inlinemetanote{writes: $S_2 :$}. Do you accept this drawing as a function: $1$ goes to $1$, $2$ goes to $2$ \inlinemetanote{draws a diagram mapping $1 \to 1$ and $2 \to 2$}? That's one element. Another element is $1$ goes to $2$ and $2$ goes to $1$ \inlinemetanote{draws a diagram mapping $1 \to 2$ and $2 \to 1$}. And those are pretty much the only two bijections, right?
\end{speech}

\begin{content}
\begin{exercise}[{Group Axioms for $S_n$}]\label[exercise]{ex:sn-group-axioms}
Verify that $(S_n, \circ, \operatorname{Id})$ satisfies the group axioms: closure under composition, associativity of $\circ$, the unit property for $\operatorname{Id}$, and existence of a bijective inverse $f^{-1} \in S_n$ for every $f \in S_n$.
\end{exercise}

\begin{proof}[Proof (beyond the lecture)]
The lecture leaves the verification of the group axioms for $S_n$ to the students (\qt{I leave all of this to you}). Let $f, g, h \in S_n$.

\textbf{Closure:} The composition $f \circ g \colon [n] \to [n]$ of two bijections is again a bijection: for $a, b \in [n]$, if $f(g(a)) = f(g(b))$, injectivity of $f$ gives $g(a) = g(b)$ and injectivity of $g$ gives $a = b$; for any $c \in [n]$, surjectivity of $f$ gives $b \in [n]$ with $f(b) = c$ and surjectivity of $g$ gives $a \in [n]$ with $g(a) = b$, so $(f \circ g)(a) = c$.

\textbf{Associativity:} For all $k \in [n]$,
\[
\bigl((f \circ g) \circ h\bigr)(k) = (f \circ g)(h(k)) = f\bigl(g(h(k))\bigr) = f\bigl((g \circ h)(k)\bigr) = \bigl(f \circ (g \circ h)\bigr)(k).
\]

\textbf{Unit property:} Clearly $\operatorname{Id} \in S_n$, and $(\operatorname{Id} \circ f)(k) = f(k) = (f \circ \operatorname{Id})(k)$ for all $k \in [n]$.

\textbf{Inverses:} Since $f$ is bijective, its set-theoretic inverse $f^{-1} \colon [n] \to [n]$ is also bijective and satisfies $f \circ f^{-1} = f^{-1} \circ f = \operatorname{Id}$.
\end{proof}

\begin{examples}[{Representing Elements of $S_1$ and $S_2$ via Arrow Diagrams}]\label[examples]{exm:s1-s2-diagrams}
For small $n$, bijections on $[n]$ can be drawn explicitly as arrow diagrams:
\definecolor{figBlue}{HTML}{07608F} % RGB (7, 96, 143): petrol blue (MidnightBlue)
\definecolor{figOrange}{HTML}{CC7E2B} % RGB (204, 126, 43): ochre orange (BurntOrange!80!black)
\begin{hidden}
- Canvas: no scale, cm; three arrow diagrams side by side: S_1 (unique map), S_2 (identity), S_2 (transposition).
  Panel centres at x = 0.9, 5.1, 9.3 (shifted by 4.2 cm each); picture spans x in [-0.8, 11.0], y in [-0.9, 1.5]
- Order: (1) arrows with shorten >=3pt, shorten <=3pt so arrowheads clear the dots; (2) dots; (3) all text last
- Panel 1 (S_1): dots at (0, 1.2) and (1.8, 1.2); labels "1" anchor east at (0, 1.2) (x in [-0.35, -0.15]) and anchor
  west at (1.8, 1.2) (x in [1.95, 2.15]), inner sep 0.15cm
- Panel 2 (S_2 identity): dots at (4.2, 1.2), (4.2, 0.2), (6.0, 1.2), (6.0, 0.2); labels "1", "2" anchor east at x = 4.2
  and anchor west at x = 6.0; vertical dot spacing 1.0 cm > 0.45 + 0.15 = 0.60 cm
- Panel 3 (S_2 swap): dots at (8.4, 1.2), (8.4, 0.2), (10.2, 1.2), (10.2, 0.2); crossed arrows in figOrange
- Gaps: Panel 1 right label ends at 2.15, Panel 2 left label starts at 3.85 (gap 1.70); Panel 2 right ends at 6.35,
  Panel 3 left starts at 8.05 (gap 1.70). Captions (text width 3.4cm, align=flush center) anchor north at y = -0.35:
  centres 4.2 cm apart -> 0.8 cm gap between neighbouring caption boxes.
\end{hidden}
\begin{center}
\begin{tikzpicture}[>=Stealth]
    % (1) Arrows
    \draw[->, thick, figBlue, shorten >=3pt, shorten <=3pt] (0, 1.2) -- (1.8, 1.2);
    \draw[->, thick, figBlue, shorten >=3pt, shorten <=3pt] (4.2, 1.2) -- (6.0, 1.2);
    \draw[->, thick, figBlue, shorten >=3pt, shorten <=3pt] (4.2, 0.2) -- (6.0, 0.2);
    \draw[->, thick, figOrange, shorten >=3pt, shorten <=3pt] (8.4, 1.2) -- (10.2, 0.2);
    \draw[->, thick, figOrange, shorten >=3pt, shorten <=3pt] (8.4, 0.2) -- (10.2, 1.2);

    % (2) Dots
    \fill (0, 1.2) circle (2pt) (1.8, 1.2) circle (2pt);
    \fill (4.2, 1.2) circle (2pt) (4.2, 0.2) circle (2pt) (6.0, 1.2) circle (2pt) (6.0, 0.2) circle (2pt);
    \fill (8.4, 1.2) circle (2pt) (8.4, 0.2) circle (2pt) (10.2, 1.2) circle (2pt) (10.2, 0.2) circle (2pt);

    % (3) Text last
    \node[anchor=east, inner sep=0.15cm] at (0, 1.2) {$1$};
    \node[anchor=west, inner sep=0.15cm] at (1.8, 1.2) {$1$};
    \node[anchor=east, inner sep=0.15cm] at (4.2, 1.2) {$1$};
    \node[anchor=east, inner sep=0.15cm] at (4.2, 0.2) {$2$};
    \node[anchor=west, inner sep=0.15cm] at (6.0, 1.2) {$1$};
    \node[anchor=west, inner sep=0.15cm] at (6.0, 0.2) {$2$};
    \node[anchor=east, inner sep=0.15cm] at (8.4, 1.2) {$1$};
    \node[anchor=east, inner sep=0.15cm] at (8.4, 0.2) {$2$};
    \node[anchor=west, inner sep=0.15cm] at (10.2, 1.2) {$1$};
    \node[anchor=west, inner sep=0.15cm] at (10.2, 0.2) {$2$};
    \node[anchor=north, text width=3.4cm, align=flush center, inner sep=0.15cm] at (0.9, -0.35) {$S_1$: unique map};
    \node[anchor=north, text width=3.4cm, align=flush center, inner sep=0.15cm] at (5.1, -0.35) {$S_2$: identity};
    \node[anchor=north, text width=3.4cm, align=flush center, inner sep=0.15cm] at (9.3, -0.35) {$S_2$: swap};
\end{tikzpicture}
\end{center}
\end{examples}
\end{content}

\begin{speech}[01:04:54 - 01:05:56]
But this is kind of unfortunate --- if we're going to have to draw a picture for every element, that's kind of unfortunate. Like $S_3$, you know, things get more complicated: you could have $1, 2, 3$ --- you could have $1$ goes to $2$, $2$ goes to $3$, $3$ goes to $1$; that's a nice element \inlinemetanote{writes \qt{$S_3:$} and draws the diagram mapping $1 \to 2$, $2 \to 3$, $3 \to 1$}. But you can imagine, this is not a very good way of naming the function, by just giving its full... somehow it takes too much chalk, let's say. So, that's one element.

I should say that it should be clear that $S_n$ is a finite group, and its order --- I mean, the standard notation for order is putting these straight lines --- is $n$ factorial \inlinemetanote{writes \qt{$S_n \text{ is finite group}$} and \qt{$|S_n| = n!$}}. That's how many bijections there are, and that's a standard formula, because you can take the first one to any $n$, and then that position's occupied, the second one has to go to $n - 1$, and so on. That's kind of a proof.
\end{speech}

\begin{content}[{Order of $S_n$ and an Element of $S_3$}]
As $n$ grows, drawing arrow diagrams for individual elements of $S_n$ (such as the cyclic permutation in $S_3$ shown below) quickly becomes unwieldy:
\definecolor{figBlue}{HTML}{07608F} % RGB (7, 96, 143): petrol blue (MidnightBlue)
\begin{hidden}
- Canvas: no scale, cm; bipartite diagram for the 3-cycle 1 -> 2 -> 3 -> 1 in S_3.
  Left dots at x = 0, right dots at x = 2.4, y = 2.0, 1.0, 0.0; picture spans x in [-0.4, 2.8], y in [-0.4, 2.3]
- Order: (1) arrows with shorten >=3pt, shorten <=3pt; (2) dots; (3) text last
- Labels: "1", "2", "3" anchor east at (0, y) and anchor west at (2.4, y) with inner sep 0.15cm.
- Collisions: vertical spacing 1.0 cm clears 0.45 cm label heights; arrows run strictly inside x in [0, 2.4].
\end{hidden}
\begin{center}
\begin{tikzpicture}[>=Stealth]
    % (1) Arrows: 1 -> 2, 2 -> 3, 3 -> 1
    \draw[->, thick, figBlue, shorten >=3pt, shorten <=3pt] (0, 2.0) -- (2.4, 1.0);
    \draw[->, thick, figBlue, shorten >=3pt, shorten <=3pt] (0, 1.0) -- (2.4, 0.0);
    \draw[->, thick, figBlue, shorten >=3pt, shorten <=3pt] (0, 0.0) -- (2.4, 2.0);

    % (2) Dots
    \fill (0, 2.0) circle (2pt) (0, 1.0) circle (2pt) (0, 0.0) circle (2pt);
    \fill (2.4, 2.0) circle (2pt) (2.4, 1.0) circle (2pt) (2.4, 0.0) circle (2pt);

    % (3) Text last
    \node[anchor=east, inner sep=0.15cm] at (0, 2.0) {$1$};
    \node[anchor=east, inner sep=0.15cm] at (0, 1.0) {$2$};
    \node[anchor=east, inner sep=0.15cm] at (0, 0.0) {$3$};
    \node[anchor=west, inner sep=0.15cm] at (2.4, 2.0) {$1$};
    \node[anchor=west, inner sep=0.15cm] at (2.4, 1.0) {$2$};
    \node[anchor=west, inner sep=0.15cm] at (2.4, 0.0) {$3$};
\end{tikzpicture}
\end{center}

\begin{proposition}[{Order of the Symmetric Group $S_n$}]\label[proposition]{prop:order-symmetric-group}
The symmetric group $S_n$ is a finite group of order
\[
|S_n| = n!.
\]
\end{proposition}

\begin{proof}
To construct a bijection $f \colon [n] \to [n]$, there are $n$ choices for $f(1)$. Once $f(1)$ is chosen, that image is occupied, leaving $n - 1$ choices for $f(2)$. Continuing inductively, there are $n - k + 1$ choices for $f(k)$, down to $1$ choice for $f(n)$. Hence, there are $n(n-1)\cdots 1 = n!$ elements in $S_n$.
\end{proof}
\end{content}

\begin{speech}[01:06:04 - 01:06:39]
Okay, so, far, the question is: how are we going to talk about these elements of the symmetric group? And there's a standard way, which is about the \newterm{cycle decomposition}, so I'll try to write it, and it's really useful notation.

In some sense it's again about names. We need to find some nice way to name these elements that will allow us to manipulate them both in our head and on the board and on paper. Logically, one doesn't have to do this --- logically, you can view those diagrams as the data type. \inlinemetanote{erases the lower chalkboard}
\end{speech}

\begin{hidden}
% timestamp: 01:07:00
% topic: Definition and order of the symmetric group S_n
% board_state: not:n-element-set, def:symmetric-group, rem:composition-order, ex:sn-group-axioms, exm:s1-s2-diagrams, prop:order-symmetric-group
% next_goal: Introduce cycle decomposition notation, list all elements of S_3, compute inverses and products
% open_loops: none
\end{hidden}

\subsection{Cycle Decomposition Notation}

\begin{speech}[01:07:00 - 01:08:56]
But there is a standard way to do it, which is the cycle decomposition, and once you understand it --- maybe you already understand it --- you'll find this a very nice way to do it \inlinemetanote{writes \qt{$\text{How to name elmts. of } S_n$} and writes and underlines \qt{$\text{Cycle decomp. notation}$}}.

And how this works: let me do an example. I'll do it in a couple examples, and then you should be able to see. So, if I take the $3$ case, let's do this one: $1$ goes to $2$, $2$ goes to $1$, $3$ goes to $3$ \inlinemetanote{draws the arrow diagram mapping $1 \to 2$, $2 \to 1$, $3 \to 3$}. So, this is one of the bijections: I flip $1$ and $2$, and $3$ stays the same. Okay?

And the way to represent it in the cycle decomposition is I start by --- I open a parenthesis, and I put $1$ here \inlinemetanote{writes: $(1$}. And then next to $1$, I write where $1$ goes. So, where does $1$ go? It goes to $2$ \inlinemetanote{writes: $2$}. Okay? And then I look, and where does $2$ go? Well, $2$ goes back to $1$, so then that closes the cycle \inlinemetanote{closes the parenthesis, giving \qt{$(1\,2)$}}.

And then I take the first thing that's not here, which is $3$, so I open a parenthesis and I see where $3$ goes, but $3$ goes back to itself \inlinemetanote{writes: $(3)$}. So, the cycle decomposition of this would be that \inlinemetanote{points at $(1\,2)(3)$}. All right?

Another example: the one I had there, this one \inlinemetanote{draws the arrow diagram mapping $1 \to 2$, $2 \to 3$, $3 \to 1$}. So, again, I start, I open the parenthesis and there's $1$; $1$ goes to $2$ again; now $2$ goes to $3$; and $3$ goes back to $1$, I close the cycle \inlinemetanote{writes: $(1\,2\,3)$}. So, that's a...
\end{speech}

\begin{content}[Cycle Decomposition Examples in \texorpdfstring{$S_3$}{S\_3}]
To name elements of $S_n$ compactly, we write each permutation in \newterm{cycle decomposition notation}. Starting from an element, we open a parenthesis, write each successive image beside its preimage, and close the parenthesis once the chain returns to the starting element; we then repeat from the smallest element not yet listed.

\begin{definition}[Cycle]\label[definition]{def:cycle}
A permutation $\sigma \in S_n$ is called a \newterm{$k$-cycle} (or cycle of length $k$), denoted $(i_1 \, i_2 \, \dots \, i_k)$, if
\[
\sigma(i_1) = i_2, \quad \sigma(i_2) = i_3, \quad \dots, \quad \sigma(i_{k-1}) = i_k, \quad \sigma(i_k) = i_1,
\]
and $\sigma(j) = j$ for all $j \in [n] \setminus \{i_1, \dots, i_k\}$.
\end{definition}

\definecolor{figBlue}{HTML}{07608F} % RGB (7, 96, 143): petrol blue (MidnightBlue)
\definecolor{figGray}{HTML}{666665} % RGB (102, 102, 101): dark gray (gray!80!black)
\begin{hidden}
- Canvas: no scale, cm; two bipartite arrow diagrams in S_3 side by side: (1 2)(3) on the left (x = 0 to 2.2),
  (1 2 3) on the right (x = 5.6 to 7.8); y dots at 2.0, 1.0, 0.0; picture spans x in [-0.4, 8.2], y in [-1.1, 2.3]
- Order: (1) arrows with shorten >=3pt, shorten <=3pt; (2) dots; (3) text last (1, 2, 3 labels and cycle captions)
- Left diagram (1 2)(3): 1 -> 2 and 2 -> 1 in figBlue, 3 -> 3 in figGray; caption "(1 2)(3)" anchor north at (1.1, -0.45)
- Right diagram (1 2 3): 1 -> 2, 2 -> 3, 3 -> 1 in figBlue; caption "(1 2 3)" anchor north at (6.7, -0.45)
- Gaps: left panel right labels end at x = 2.55, right panel left labels start at x = 5.25 (gap 2.70 cm);
  vertical spacing between rows 1.0 cm clears 0.45 cm labels.
\end{hidden}
\begin{center}
\begin{tikzpicture}[>=Stealth]
    % (1) Arrows: Left = (1 2)(3)
    \draw[->, thick, figBlue, shorten >=3pt, shorten <=3pt] (0, 2.0) -- (2.2, 1.0);
    \draw[->, thick, figBlue, shorten >=3pt, shorten <=3pt] (0, 1.0) -- (2.2, 2.0);
    \draw[->, thick, figGray, shorten >=3pt, shorten <=3pt] (0, 0.0) -- (2.2, 0.0);

    % Right = (1 2 3)
    \draw[->, thick, figBlue, shorten >=3pt, shorten <=3pt] (5.6, 2.0) -- (7.8, 1.0);
    \draw[->, thick, figBlue, shorten >=3pt, shorten <=3pt] (5.6, 1.0) -- (7.8, 0.0);
    \draw[->, thick, figBlue, shorten >=3pt, shorten <=3pt] (5.6, 0.0) -- (7.8, 2.0);

    % (2) Dots
    \fill (0, 2.0) circle (2pt) (0, 1.0) circle (2pt) (0, 0.0) circle (2pt);
    \fill (2.2, 2.0) circle (2pt) (2.2, 1.0) circle (2pt) (2.2, 0.0) circle (2pt);
    \fill (5.6, 2.0) circle (2pt) (5.6, 1.0) circle (2pt) (5.6, 0.0) circle (2pt);
    \fill (7.8, 2.0) circle (2pt) (7.8, 1.0) circle (2pt) (7.8, 0.0) circle (2pt);

    % (3) Text last
    \node[anchor=east, inner sep=0.15cm] at (0, 2.0) {$1$};
    \node[anchor=east, inner sep=0.15cm] at (0, 1.0) {$2$};
    \node[anchor=east, inner sep=0.15cm] at (0, 0.0) {$3$};
    \node[anchor=west, inner sep=0.15cm] at (2.2, 2.0) {$1$};
    \node[anchor=west, inner sep=0.15cm] at (2.2, 1.0) {$2$};
    \node[anchor=west, inner sep=0.15cm] at (2.2, 0.0) {$3$};

    \node[anchor=east, inner sep=0.15cm] at (5.6, 2.0) {$1$};
    \node[anchor=east, inner sep=0.15cm] at (5.6, 1.0) {$2$};
    \node[anchor=east, inner sep=0.15cm] at (5.6, 0.0) {$3$};
    \node[anchor=west, inner sep=0.15cm] at (7.8, 2.0) {$1$};
    \node[anchor=west, inner sep=0.15cm] at (7.8, 1.0) {$2$};
    \node[anchor=west, inner sep=0.15cm] at (7.8, 0.0) {$3$};

    \node[anchor=north, inner sep=0.15cm] at (1.1, -0.45) {$(1\,2)(3)$};
    \node[anchor=north, inner sep=0.15cm] at (6.7, -0.45) {$(1\,2\,3)$};
\end{tikzpicture}
\end{center}
\end{content}

\begin{speech}[01:08:56 - 01:10:55]
So, if I write down now what is $S_3$, I can write down all the cycle decompositions \inlinemetanote{writes: $S_3 =$}. The first one is the trivial one; that's where no one goes anywhere \inlinemetanote{writes: $(1)(2)(3)$}. Then there's these ones: $(1\,2)(3)$, $(1\,3)(2)$, $(1)(2\,3)$ \inlinemetanote{writes \qt{$(1\,2)(3)$}, \qt{$(1\,3)(2)$}, and \qt{$(1)(2\,3)$}}. And that's already $4$, I need $2$ more: there's $(1\,2\,3)$ and $(1\,3\,2)$ \inlinemetanote{writes \qt{$(1\,2\,3)$} and \qt{$(1\,3\,2)$} below}. So, here I've listed all $6$ elements.

Do you understand this notation? Like, if I pick this \inlinemetanote{points at $(1\,2\,3)$} --- these cycles just tell you where everything goes. Like, if I take this one \inlinemetanote{points at $(1\,2\,3)$} and I want to know where $2$ goes, I find $2$, and then it goes to the next one ($3$). And $3$, if I want to know $3$ --- whenever the cycle ends, it means you go to the first one ($1$). And if you're by yourself \inlinemetanote{points at $(1)$ in $(1)(2)(3)$}, it means you don't go anywhere: you're a fixed point.

So, are you satisfied with this? You have to play with it. But it's much more efficient, because I've already written for you the $6$ elements in $S_3$.

Traditionally --- I mean, the custom is you don't write these fixed points, because they're kind of a waste of chalk also. So, the standard is that if they don't appear, they're fixed. So, the standard way I would write this would be: that, I would just call it $e$, the neutral element \inlinemetanote{writes \qt{Standard} underlined and \qt{$e$} below it}; and this I would say $(1\,2)$, $(1\,3)$, $(2\,3)$, $(1\,2\,3)$, $(1\,3\,2)$ \inlinemetanote{writes \qt{$(1\,2)$}, \qt{$(1\,3)$}, \qt{$(2\,3)$}, \qt{$(1\,2\,3)$}, and \qt{$(1\,3\,2)$}}. That would be the standard: if they're not there, it means nothing's happening to them. That's kind of the standard way to represent it.

And, you know, this way I picked some way to break symmetry, but, you know, you could also write this $(2\,3)$ as $(3\,2)$, because $2$ goes to $3$ and $3$ goes to $2$, and here $3$ goes to $2$, $2$ goes to $3$ \inlinemetanote{writes \qt{$(3\,2)$} to the left of $(2\,3)$}. You can also write it like that if you want, but this \inlinemetanote{points at the Standard list} is a canonical way of doing it. All right.
\end{speech}

\begin{content}
\begin{examples}[{All Six Elements of $S_3$ in Full and Standard Cycle Notation}]\label[examples]{exm:s3-cycles}
Listing every cycle explicitly (including $1$-cycles for fixed points) versus the standard shorthand (omitting $1$-cycles and denoting the identity by $e$):
\begin{center}
\begin{tabular}{l l}
    \toprule
    \textbf{Full Cycle Decomposition} & \textbf{Standard Notation} \\
    \midrule
    $(1)(2)(3)$   & $e$ \\
    $(1\,2)(3)$   & $(1\,2)$ \\
    $(1\,3)(2)$   & $(1\,3)$ \\
    $(1)(2\,3)$   & $(2\,3)$ \\
    $(1\,2\,3)$   & $(1\,2\,3)$ \\
    $(1\,3\,2)$   & $(1\,3\,2)$ \\
    \bottomrule
\end{tabular}
\end{center}
Note that a cycle can also be written starting from a different entry --- for instance, $(2\,3) = (3\,2)$ --- whereas writing the smallest element of each cycle first gives a canonical way of representing it.
\end{examples}

\begin{notation}[Suppression of $1$-Cycles]\label[notation]{not:omitting-1-cycles}
In standard cycle notation, $1$-cycles (fixed points) are omitted: the identity permutation $(1)(2)\cdots(n)$ is denoted $e$ (or $\operatorname{Id}$), and a cycle permutation such as $(i\,j)(k)$ is written compactly as $(i\,j)$ with fixed points left implicit.
\end{notation}
\end{content}

\begin{aiNote}[Notation: Transpositions]
% [PERSISTENT]
Many books call a $2$-cycle a \emph{transposition}. The lecturer mostly says \qt{two-cycle} (using \qt{transposition} once in \cref{prop:im-phi-hat}).
\end{aiNote}

\begin{aiNote}[Beyond the Lecture: The Inverse of a Cycle]
% [PERSISTENT, ADDITION]
The inverse of any cycle $(i_1 \, i_2 \, \dots \, i_k)$ is obtained by reversing the order of its entries:
\[
(i_1 \, i_2 \, \dots \, i_k)^{-1} = (i_k \, i_{k-1} \, \dots \, i_1).
\]
For example, $(1\,2\,3)^{-1} = (3\,2\,1) = (1\,3\,2)$.
\end{aiNote}

\begin{speech}[01:10:55 - 01:12:53]
So, that works of course for any $n$, and the rule is: if you have this kind of function, it tells you the cycle decomposition; if you have the cycle decomposition, it tells you the function. But it's much more efficient.

And so let us try to... in this notation, what is the inverse map \inlinemetanote{writes and underlines \qt{Inv}, with an arrow from \qt{Standard} to \qt{Inv}}?

Of course, the inverse of $e$ is, of course, $e$, right \inlinemetanote{draws an arrow from $e$ to $e$}? Neutral element. What's the inverse of this element \inlinemetanote{points at $(1\,2)$}? \inlinemetanote{inaudible answer from a student} Yeah. Itself \inlinemetanote{draws an arrow from $(1\,2)$ to $(1\,2)$}.

If I apply it twice --- so, now I have to tell you how to multiply in the cycle notation. So, here I take $(1\,2)$, that's $1$ times --- I can just put the cycles together; that's it, I just put them together \inlinemetanote{writes: $(1\,2)(1\,2) =$}. And now I have to figure out what cycle to write here.

So, then I have to find out where $1$ --- I open it up, I find out where $1$ goes \inlinemetanote{writes: $($}: $1$ goes to $2$, and then I go to the next one, now $2$ goes back to $1$, so $1$ stays by itself \inlinemetanote{writes: $1)($}. And I start here, $2$: $2$ goes to $1$, $1$ goes back to $2$, so $2$ goes to $2$ \inlinemetanote{writes: $2)$}. That's the identity. Also you can just see it in functions. But there's just rules.

Of course, these are all the same, so this is the inverse; that's the inverse \inlinemetanote{draws arrows from $(1\,3)$ to $(1\,3)$ and from $(2\,3)$ to $(2\,3)$}. What's the inverse of this person \inlinemetanote{points at $(1\,2\,3)$}? \inlinemetanote{inaudible answer from a student} Yeah, the other one. So, they flip around here: $(1\,3\,2)$ and $(1\,2\,3)$ \inlinemetanote{draws arrows $(1\,2\,3) \to (1\,3\,2)$ and $(1\,3\,2) \to (1\,2\,3)$}. So, that's the inverse map.

These are more complicated: $(1\,2\,3)$, let us say, $(1\,3\,2)$ \inlinemetanote{writes: $(1\,2\,3)(1\,3\,2) =$}. So, let's do this one slowly, just for fun --- I don't know how fun it is, but anyway: $1$ --- so, to figure out where $1$ goes, you start with $1$ in the first cycle you see (on the right), where $1$ is, and this cycle tells you $1$ goes to $3$. And then you keep going left: ah, you still see a $3$, so $3$ is going somewhere; $3$ goes back to $1$ \inlinemetanote{writes \qt{$1 \to 3 \to 1$} below}. That's kind of nice. So, that's $(1)$, and you'll see $(2)$ and $(3)$ \inlinemetanote{writes \qt{$(1)$} and \qt{$(2)(3)$}}.
\end{speech}

\begin{content}[{Inverses in $S_3$ and Multiplying Cycles}]
To multiply two permutations in cycle notation, we concatenate them and trace each element from right to left (matching function composition $f \circ g$). Applying this gives the inverse map $\operatorname{Inv} \colon S_3 \to S_3$:
\begin{center}
\begin{tabular}{l c l}
    \toprule
    \textbf{Element} & & \textbf{Inverse} \\
    \midrule
    $e$          & $\longmapsto$ & $e$ \\
    $(1\,2)$     & $\longmapsto$ & $(1\,2)$ \\
    $(1\,3)$     & $\longmapsto$ & $(1\,3)$ \\
    $(2\,3)$     & $\longmapsto$ & $(2\,3)$ \\
    $(1\,2\,3)$  & $\longmapsto$ & $(1\,3\,2)$ \\
    $(1\,3\,2)$  & $\longmapsto$ & $(1\,2\,3)$ \\
    \bottomrule
\end{tabular}
\end{center}
Direct cycle multiplication verifies these inverses:
\begin{align*}
(1\,2)(1\,2) &= (1)(2) && (= e), \\
(1\,2\,3)(1\,3\,2) &= (1)(2)(3) && (= e).
\end{align*}

\begin{steps}[Tracing Elements Right-to-Left]
In $(1\,2\,3)(1\,3\,2)$:
\begin{itemize}
    \item $1 \mapsto 3$ (under $(1\,3\,2)$) $\mapsto 1$ (under $(1\,2\,3)$), closing the $1$-cycle $(1)$.
    \item $2 \mapsto 1$ (under $(1\,3\,2)$) $\mapsto 2$ (under $(1\,2\,3)$), closing $(2)$.
    \item $3 \mapsto 2$ (under $(1\,3\,2)$) $\mapsto 3$ (under $(1\,2\,3)$), closing $(3)$.
\end{itemize}
Hence, $(1\,2\,3)(1\,3\,2) = (1)(2)(3) = e$.
\end{steps}
\end{content}

\begin{speech}[01:12:53 - 01:14:18]
Okay, let's try some... That's the inverse map, but we should be able to multiply anything. Like, what is $(1\,3)$ times $(2\,3)$ \inlinemetanote{writes: $(1\,3)(2\,3) =$}?

So, to do this, as I said, just start out with $1$: $1$ doesn't do anything here (in $(2\,3)$), but it goes to $3$ (in $(1\,3)$), so over there \inlinemetanote{writes: $(1\,3$}. So, now we do $3$: $3$ goes to $2$ (in $(2\,3)$); $2$ doesn't go anywhere here (in $(1\,3)$), right? So, it stays at $2$ \inlinemetanote{writes: $2$}. $2$ goes to $3$ (in $(2\,3)$), and $3$ goes to $1$ (in $(1\,3)$), so it's that cycle \inlinemetanote{closes the parenthesis, giving \qt{$(1\,3\,2)$}}. That's the product.

And as I said, this is notation; you have to see how it works. You could, if you're suspicious about this notation, every time you see this, make one of those diagrams and put them together \inlinemetanote{points at the arrow diagrams on the left board}, but it's not a good idea; it's better to get to understand this notation. That's how people actually work with the cycles.

So, an exercise that you should do, if you've never seen it... but have you seen this kind of cycle notation in other classes? No? Some people have, some people haven't.

It's just notation, and there's not much more I can say: you have to become familiar with it yourself. \inlinemetanote{erases the lower left arrow diagram and middle-left board silently}
\end{speech}

\begin{content}[{Sample Product in $S_3$: $(1\,3)(2\,3)$}]
Computing the product $(1\,3)(2\,3)$ directly from right to left yields:
\[
(1\,3)(2\,3) = (1\,3\,2).
\]

\begin{steps}[Step-by-Step Cycle Tracing]
\begin{itemize}
    \item Starting at $1$: $(2\,3)$ fixes $1$, and $(1\,3)$ sends $1 \mapsto 3$, so $1 \mapsto 3$.
    \item Continuing from $3$: $(2\,3)$ sends $3 \mapsto 2$, and $(1\,3)$ fixes $2$, so $3 \mapsto 2$.
    \item Continuing from $2$: $(2\,3)$ sends $2 \mapsto 3$, and $(1\,3)$ sends $3 \mapsto 1$, closing the cycle $(1\,3\,2)$.
\end{itemize}
\end{steps}
\end{content}

\begin{hidden}
% timestamp: 01:14:18
% topic: Cycle decomposition notation, full list of S_3, inverses, and sample cycle products
% board_state: def:cycle, exm:s3-cycles, not:omitting-1-cycles
% next_goal: Compute (2 3)(1 3) to illustrate non-commutativity, assign S_3 multiplication table exercise, then study finite groups
% open_loops: none
\end{hidden}

\begin{speech}[01:14:45 - 01:15:45]
Oh, I should do one more example while I'm there, by the way. Let's do one more. What if I do it in the different way: $(2\,3)(1\,3)$ \inlinemetanote{writes: $(2\,3)(1\,3) =$}?

Then... so now you tell me: $1$, where does $1$ go? Someone should say $2$. Yeah, $1$ goes to $2$, because you go: $1$ goes to $3$ (in $(1\,3)$), then $3$ gets bumped to $2$ (in $(2\,3)$) \inlinemetanote{writes: $(1\,2$}.

And where does $2$ go? Well, $2$ doesn't see anything here (in $(1\,3)$), but then gets bumped to $3$ (in $(2\,3)$), and then $3$ gets bumped to $1$ here (in $(1\,3)$)... so that's it \inlinemetanote{writes \qt{$3$} and closes the parenthesis, giving \qt{$(1\,2\,3)$}}. So, actually, this is already showing that it's not commutative, because you got the two different elements when you do them in different order \inlinemetanote{points at $(1\,3)(2\,3) = (1\,3\,2)$ and $(2\,3)(1\,3) = (1\,2\,3)$}.

So, what I advocate --- what I was trying to say --- is I advocate, if you are learning this cycle notation, to go home and write the multiplication table of $S_3$.
\end{speech}

\begin{content}[{Non-Commutativity of $S_3$: Comparing $(1\,3)(2\,3)$ and $(2\,3)(1\,3)$}]
\begin{observation}[Non-Commutativity in $S_3$]\label[observation]{obs:s3-non-commutative}
Reversing the order of factors in the product above gives:
\[
(2\,3)(1\,3) = (1\,2\,3).
\]
Since $(1\,3)(2\,3) = (1\,3\,2) \neq (1\,2\,3) = (2\,3)(1\,3)$, this example explicitly shows that $S_3$ is non-commutative.
\end{observation}
\end{content}

\begin{aiNote}[Beyond the Lecture: The Smallest Non-Abelian Group]
% [PERSISTENT, ADDITION]
$S_3$, of order $6$, is the smallest non-abelian group.
\end{aiNote}

\begin{speech}[01:15:45 - 01:17:00]
So, $S_3$ --- it's not that many; I mean, most of them we know, we've already done. $S_3$ --- the multiplication table, of course, is this table where I write the elements of $S_3$ here, and I write dot, and then I write them again here \inlinemetanote{draws a table header above the $S_3$ column with \qt{$\cdot$} and \qt{$(1)(2)(3) \ \dots \ (1\,3\,2)$}}. So, there's $36$ elements and you just fill it out --- that's a good exercise. Most of them we've already done, actually. After you do that, you'll understand.

You can do it for $S_4$ too, but it has more elements: that has $576$ elements, so you might... \inlinemetanote{erases the lower chalkboard} take the afternoon off and do that.

But one way or the other, you should make sure you understand what this symmetric group is, how to think about its elements, how to compose --- it's a kind of important part of this class. And whatever direction your future mathematical life will go, it will have the symmetric group. It's like one of the really basic objects in mathematics. You cannot get away from the symmetric group.
\end{speech}

\begin{content}
\begin{exercise}[{Multiplication Table of $S_3$}]\label[exercise]{ex:s3-cayley-table}
Fill out the $6 \times 6$ ($36$-entry) multiplication table of $S_3$ (whose header was sketched on the board as shown below, with the column header retaining the full $1$-cycle form $(1)(2)(3)$ for $e$):
\begin{center}
\begin{tabular}{c | c c c c c c}
    $\cdot$ & $(1)(2)(3)$ & $(1\,2)$ & $(1\,3)$ & $(2\,3)$ & $(1\,2\,3)$ & $(1\,3\,2)$ \\
    \hline
    $e$         & & & & & & \\
    $(1\,2)$    & & & & & & \\
    $(1\,3)$    & & & & & & \\
    $(2\,3)$    & & & & & & \\
    $(1\,2\,3)$ & & & & & & \\
    $(1\,3\,2)$ & & & & & &
\end{tabular}
\end{center}
(In contrast, the table for $S_4$ has $24 \times 24 = 576$ entries.)
\end{exercise}

\begin{proof}[Solution to \cref{ex:s3-cayley-table} (beyond the lecture)]
The lecture leaves filling out the $6 \times 6$ multiplication table of $S_3$ as an exercise. Using right-to-left composition ($f \cdot g$ means apply $g$ first, then $f$), where the row labels give $f$ and the column labels give $g$:
\begin{center}
\begin{tabular}{c | c c c c c c}
    \toprule
    \diagbox{$f$}{$g$} & $e$ & $(1\,2)$ & $(1\,3)$ & $(2\,3)$ & $(1\,2\,3)$ & $(1\,3\,2)$ \\
    \midrule
    $e$         & $e$         & $(1\,2)$    & $(1\,3)$    & $(2\,3)$    & $(1\,2\,3)$ & $(1\,3\,2)$ \\
    $(1\,2)$    & $(1\,2)$    & $e$         & $(1\,3\,2)$ & $(1\,2\,3)$ & $(2\,3)$    & $(1\,3)$ \\
    $(1\,3)$    & $(1\,3)$    & $(1\,2\,3)$ & $e$         & $(1\,3\,2)$ & $(1\,2)$    & $(2\,3)$ \\
    $(2\,3)$    & $(2\,3)$    & $(1\,3\,2)$ & $(1\,2\,3)$ & $e$         & $(1\,3)$    & $(1\,2)$ \\
    $(1\,2\,3)$ & $(1\,2\,3)$ & $(1\,3)$    & $(2\,3)$    & $(1\,2)$    & $(1\,3\,2)$ & $e$ \\
    $(1\,3\,2)$ & $(1\,3\,2)$ & $(2\,3)$    & $(1\,2)$    & $(1\,3)$    & $e$         & $(1\,2\,3)$ \\
    \bottomrule
\end{tabular}
\end{center}
\end{proof}
\end{content}

\begin{aiNote}[Beyond the Lecture: Cayley's Theorem]
% [PERSISTENT, ADDITION]
As the lecturer remarks about the ubiquity of $S_n$, Cayley's Theorem (\cref{thm:cayley}, \cref{thm:cayley-embedding}) makes this precise: every finite group of order $n$ is isomorphic to a subgroup of $S_n$. Thus, permutations capture the complete structural diversity of all finite group operations.
\end{aiNote}

\begin{hidden}
% timestamp: 01:17:00
% topic: Non-commutativity of S_3 and multiplication table exercise
% board_state: obs:s3-non-commutative, ex:s3-cayley-table
% next_goal: Study finite groups, define abelian groups, prove every element of a finite group has finite order (x^m = e)
% open_loops: none
\end{hidden}

\subsection{Powers and Orders of Elements in Finite Groups}

\begin{speech}[01:17:00 - 01:18:13]
Okay, so: let $G$ be a finite group \inlinemetanote{writes: $\text{Let } G \text{ be a finite group}$}. So, this will be the assumption for many statements now for the next couple weeks.

And always by this, I, you know, talk about this $(G, \cdot, e)$ \inlinemetanote{writes: $(G, \cdot, e)$}, and I do not assume it's commutative \inlinemetanote{writes \qt{I don't assume Commutative!} and underlines \qt{not}}. I never assume it. I mean, sorry --- I never silently assume it. Sometimes I'll say \qt{let $G$ be a commutative group}, but I will never silently assume it. If I just say a group, it means that I am considering in my head a group possibly non-commutative.

I should say --- and I made this definition before --- that a commutative group: when $x \cdot y = y \cdot x$ for all $x, y$ \inlinemetanote{writes: $\text{When } x \cdot y = y \cdot x \;\; \forall x, y$}, then we say $G$ is \newterm{abelian} \inlinemetanote{writes \qt{then $G$ is Abelian} and underlines \qt{Abelian}}. That's the standard. A commutative group --- the definition, it's abelian, that's what.

So, I will... I said, if I say $G$ is a group, you should not assume it's abelian.
\end{speech}

\begin{content}[Standing Assumption: Finite Groups and Abelian Groups]
Unless explicitly stated otherwise, let $G = (G, \cdot, e)$ be a finite group; commutativity is never silently assumed.

\begin{definition}[Abelian Group]\label[definition]{def:abelian-group}
A group $(G, \cdot, e)$ is called \newterm{abelian} (or commutative) if
\[
x \cdot y = y \cdot x \qquad \forall x, y \in G.
\]
\end{definition}
\end{content}

\begin{speech}[01:18:13 - 01:20:22]
So, we take $G$ is a finite group, and let $x$ be an element of $G$ \inlinemetanote{writes: $\text{Let } x \in G$}.

All right, so the first thing I can do is consider this: $x^1$ --- $x$ to the power $1$, which is just $x$ --- $x^2, x^3$, so on \inlinemetanote{writes: (i) $x, x^2, x^3, \dots$}. And what I mean by this, which is what everybody means: $x^2$ is just $x \cdot x$, and $x^3$ is $x \cdot x \cdot x$, so on \inlinemetanote{writes \qt{$x^2 = x \cdot x$} and \qt{$x^3 = x \cdot x \cdot x$}}, okay? That's what this notation means: just apply the group operation to itself.

All right, and I just consider all of these: $x^4, x^5$, and eventually there'll be some $x^{17}$, and eventually $x^{691}$, and so on \inlinemetanote{writes: $x^4, x^5, \dots, x^{17}, \dots, x^{691}, \dots$}. And you'll see that there are infinitely many elements in this list, but $G$ is a finite group. So, two of them have to be... there must be that some of these are coincident; these cannot all be distinct \inlinemetanote{writes \qt{Can not all be distinct} with a brief note underneath regarding $G$}. They can't all be distinct, because $G$ is finite and this list is infinite. I hope you view that as a solid argument. Sometimes this is called something like the pigeonhole principle, or something like that; that seems like a fancy word for it.

So, if they're not distinct, that means there exists some $k$ less than $\ell$ such that $x^k$ is equal to $x^\ell$ \inlinemetanote{writes: $\exists\, k < \ell \text{ such that } x^k = x^\ell$}. That's what it means that they're not all distinct: if they're not all distinct, there must be at least $2$ that are the same, and then I just name them $k$ and $\ell$ with $k$ less than $\ell$. Okay?

And now, by cancellation, I can cancel all the powers of $k$ --- $k$ is less than $\ell$ ---, so by cancellation, this means the neutral element is equal to $x$ to the $\ell - k$ \inlinemetanote{writes: $\text{by cancellation} \implies e = x^{\ell - k}$}. I just keep canceling powers of $x$ until I have none here. Okay?
\end{speech}

\begin{content}[Powers of an Element in a Finite Group]
Let $x \in G$, where $G$ is a finite group. Consider the infinite list of positive powers of $x$:
\begin{equation}\label{eq:powers-list}
x, \; x^2, \; x^3, \; x^4, \; x^5, \; \dots, \; x^{17}, \; \dots, \; x^{691}, \; \dots
\end{equation}
where $x^2 = x \cdot x$, $x^3 = x \cdot x \cdot x$, and in general $x^n$ denotes the $n$-fold product of $x$ with itself.

Since all elements of the infinite list \eqref{eq:powers-list} lie in the finite group $G$, they cannot all be distinct (by the pigeonhole principle). Thus, there exist integers $1 \le k < \ell$ such that
\[
x^k = x^\ell.
\]
Canceling $k$ factors of $x$ on both sides (applying \cref{prop:cancellation-laws} $k$ times) yields
\[
e = x^{\ell - k},
\]
with $\ell - k > 0$.
\end{content}

\begin{speech}[01:20:22 - 01:21:56]
So, this is what we've proven: we have proven that for $x$ in $G$, where $G$ is finite, there exists a positive number $m$ such that $x$ to the $m$ is the neutral element \inlinemetanote{writes \qt{$\text{We have proven}$ / $\text{that for } x \in G, \; G \text{ finite}$} and writes and boxes \qt{$\exists\, m > 0, \text{ such that } x^m = e$}}.

This is what we've proven. So, for $x$ in $G$, for any element of a finite group, there's some positive number that when I raise the element to that power I get the neutral element. We've proven it.

Okay, this is not true for infinite groups, because you could take the infinite group to be the integers, and you could take our element $x$ to be $1$, and then power for the integers is adding it more and more times to itself, and you can add $1$ to itself as many times as you want; you're never going to get $0$. And the reason that argument --- the proof fails, of course --- is that for the integers, these \inlinemetanote{points at the list $x, x^2, x^3, \dots$} are all distinct. And you have to be clear, you know, when I'm talking about the integers with addition, so... It's easy to confuse yourself.
\end{speech}

\begin{content}
\begin{proposition}[Existence of a Positive Power Equal to the Identity]\label[proposition]{prop:finite-group-identity-power}
For any element $x \in G$ of a finite group $G$, there exists a positive integer $m > 0$ such that
\[
x^m = e.
\]
\end{proposition}

\begin{proof}
As derived above, the infinite sequence $x, x^2, x^3, \dots$ in the finite group $G$ must contain a repetition $x^k = x^\ell$ for some $1 \le k < \ell$. By $k$-fold cancellation (\cref{prop:cancellation-laws}), $e = x^{\ell - k}$. Setting $m := \ell - k > 0$ gives $x^m = e$.
\end{proof}

\begin{example}[Failure in Infinite Groups: \texorpdfstring{$(\mathbb{Z}, +, 0)$}{(Z, +, 0)}]\label[example]{ex:infinite-group-z-order}
\cref{prop:finite-group-identity-power} fails when $G$ is infinite. In the additive group $(\mathbb{Z}, +, 0)$, the $m$-th \qt{power} of $x = 1$ is the $m$-fold sum $m \cdot 1 = m$. Because all terms $1, 2, 3, \dots$ are pairwise distinct in $\mathbb{Z}$, there is no positive integer $m > 0$ with $m \cdot 1 = 0$.
\end{example}
\end{content}

\begin{speech}[01:21:56 - 01:23:02]
All right, so, we've proven this --- that's an extremely basic property of groups, and it's a very important one: if I have a group, a finite group, and I have an element, that element --- if I apply powers of that element --- eventually gets back to the unit, the neutral element.

And so we come to a definition \inlinemetanote{writes: $\text{Def:}$}: let $x$ be an element of $G$, $G$ finite, then the \newterm{order} of $x$ \inlinemetanote{writes \qt{$\text{Let } x \in G, \; G \text{ finite}$}, writes \qt{$\text{then } \operatorname{Ord}(x) \text{ is the}$}, and underlines \qt{$\operatorname{Ord}(x)$}} --- that's a definition --- is the minimal $m$, minimal positive $m$ such that $x^m$ is the neutral element \inlinemetanote{writes: $\min m > 0 \text{ such that } x^m = e.$}. All right?

So, once we have --- yeah, once we have \emph{some}, we can always take the minimal. And so that's a canonical number associated: anytime you have a group that's finite and you have an element, that element has a number, which is its order, and it's an extremely important aspect of that element. That's how many times you have to apply the group operation to itself to get back the neutral element.
\end{speech}

\begin{content}
\begin{definition}[Order of an Element]\label[definition]{def:element-order}
Let $G$ be a finite group and $x \in G$. The \newterm{order} of $x$, denoted $\operatorname{Ord}(x)$, is the smallest positive integer $m > 0$ such that
\[
x^m = e.
\]
\end{definition}

\begin{steps}[Why the Minimum Exists]
By \cref{prop:finite-group-identity-power}, the set $\{ m \in \mathbb{N}_{>0} \mid x^m = e \}$ is non-empty whenever $G$ is finite; by the well-ordering principle of $\mathbb{N}_{>0}$ (\cref{fact:pre-well-ordering}), it has a unique smallest element $\operatorname{Ord}(x)$.
\end{steps}
\end{content}

\begin{speech}[01:23:02 - 01:24:34]
So, let us try to think about $S_3$ \inlinemetanote{writes and underlines: $S_3$}: what is the order of $e$ \inlinemetanote{writes: $e \to$}? Well, this order of $e$ is $1$, because you don't have to --- you apply $e$ to the $1$ is $e$ (i.e., $e^1 = e$) \inlinemetanote{writes: $\operatorname{Ord}(e) = 1$}. It's kind of silly, all right, anyway.

So, what's the order of this element: $(1\,2)$ \inlinemetanote{writes: $(1\,2) \to$}? \inlinemetanote{long pause while students think; a student answers inaudibly} Yeah. $2$, because if I take --- this one is not $e$, so the order is not $1$ \inlinemetanote{writes: $(1\,2) \neq e \quad \text{ord} \neq 1$}. If I apply it to itself, I have to figure out what the square is, and the square is that --- and we calculated it, that's $e$, so the order is $2$ \inlinemetanote{writes: $(1\,2)(1\,2) = e \quad \text{ord} : 2$}.

And, of course, for all the $2$-cycles --- this thing is called a \newterm{$2$-cycle} ---, so these are all order $2$ \inlinemetanote{writes \qt{$(1\,3) \to \text{ord} = 2$} with \qt{$(2\,3)$} below it}: these are all order $2$, for all of them.

And then the other one's this one: $(1\,2\,3)$ \inlinemetanote{writes: $(1\,2\,3) \to$}. What's the order of this? Well, again, you can imagine what it is. Well, what are you going to guess this thing is? \inlinemetanote{inaudible answer from a student} $3$, yeah, $3$. Good guess \inlinemetanote{writes: $\text{ord} = 3$}.

And because you have to apply this thing $3$ times to get back to $1$: if you take $1$, when you apply it once it goes to $2$; if you apply it twice you get $3$, it's not the identity; if you do it $3$ times, you get back to $1$. So, that's that.
\end{speech}

\begin{content}
\begin{example}[{Element Orders in $S_3$}]\label[example]{ex:orders-s3}
In the symmetric group $S_3$:
\begin{itemize}
    \item $e \longrightarrow \operatorname{Ord}(e) = 1$, since $e^1 = e$.
    \item $(1\,2) \longrightarrow \operatorname{Ord}((1\,2)) = 2$, because $(1\,2) \neq e \implies \operatorname{Ord}((1\,2)) \neq 1$, whereas $(1\,2)(1\,2) = e \implies \operatorname{Ord}((1\,2)) = 2$.
    \item Every $2$-cycle (transposition) has order $2$:
    \[
    (1\,3) \longrightarrow \operatorname{Ord} = 2, \qquad (2\,3) \longrightarrow \operatorname{Ord} = 2.
    \]
    \item $(1\,2\,3) \longrightarrow \operatorname{Ord} = 3$, since $(1\,2\,3)$ sends $1 \mapsto 2$, its square sends $1 \mapsto 3$, and its cube fixes $1$ (as well as $2$ and $3$), giving $(1\,2\,3)^3 = e$.
\end{itemize}
\end{example}
\end{content}

\begin{hidden}
% timestamp: 01:24:34
% topic: Existence of finite order x^m = e, definition of order of an element, examples in S_3
% board_state: def:abelian-group, eq:powers-list, prop:finite-group-identity-power, ex:infinite-group-z-order, def:element-order, ex:orders-s3
% next_goal: Define cyclic subgroup <x>, prove its size is ord(x), state subgroup property and isomorphism with Z/ord(x)Z
% open_loops: none
\end{hidden}

\subsection{The Cyclic Subgroup Generated by an Element}

\begin{speech}[01:24:34 - 01:26:26]
Okay, and you will notice here that part of the definition... let $x$ in $G$ be an element, where $G$ is finite \inlinemetanote{writes: $x \in G, \; G \text{ finite}$}. Then I can consider this set, which is usually called this: and it's the set that contains the neutral element, $x$, and all the way up to $x$ to the order of $x$ minus $1$ \inlinemetanote{writes \qt{$\langle x \rangle =$}, \qt{$\{e, x, \dots, x^{\operatorname{Ord}(x)-1}\}$}, and adds \qt{$\text{def}$} underneath the equals sign}. Okay? I don't put order, because order is the neutral element --- that's the definition of the order, so I leave out the next one. All right?

And the claim is these are all... so, proposition --- if I'm with this definition, this is the definition, and now maybe lemma: the size of this set is exactly the order of $x$ \inlinemetanote{writes \qt{$\text{Lemma:}$}, \qt{$|\{e, x, \dots, x^{\operatorname{Ord}(x)-1}\}|$}, and \qt{$= \operatorname{Ord}(x)$}}.

And by that I mean: certainly I've written that many symbols, order of $x$, because there's $1$ to the order of $x$ minus $1$, plus the neutral element. So, there's certainly that many things I've written, but to say that this set has this order means that they're all distinct.

So, what I have to prove to prove this lemma: I have to show all these people here are distinct elements \inlinemetanote{writes \qt{I have to show all these are distinct elmts.} with an arrow pointing up to the set}. Because, you know, over here I wrote down a set of infinite... they weren't distinct \inlinemetanote{points at the earlier list $x, x^2, x^3, \dots$}.

So, why is it that when I go from the neutral up to $\operatorname{Ord}(x) - 1$, they're distinct? Why is that? Yeah.
\end{speech}

\begin{student}[Student Answer]
Because the order was defined as the minimal element... \inlinemetanote{trails off inaudibly}
\end{student}

\begin{speech}[01:26:31 - 01:27:11]
Yeah, so, look --- I mean, it could be like what happens: like this neutral element, it certainly can't be one of these \inlinemetanote{points at $x, \dots, x^{\operatorname{Ord}(x)-1}$}, because if it was like, you know, $3$ here, it can't be that, because then that would be lower than the order \inlinemetanote{writes \qt{$3$} under the dots between $x$ and $x^{\operatorname{Ord}(x)-1}$}. So, nothing here could be the neutral element.

More interesting is why can none of these two be equal \inlinemetanote{points at two powers between $x$ and $x^{\operatorname{Ord}(x)-1}$}? And again for the same reason: if these two are equal, then I could use my cancellation argument (Recall: \cref{prop:cancellation-laws}), and I would get a power of $m$ --- $m$ lower --- that's the neutral. So, you should think about that; it's an elementary argument, but it's certainly crucial. So, these things are just totally crucial aspects about group theory; we'll discuss them.
\end{speech}

\begin{content}[{The Set $\langle x \rangle$ and Its Cardinality}]
\begin{definition}[{Powers Generated by $x$}]\label[definition]{def:cyclic-subgroup}
Let $G$ be a finite group and $x \in G$. We define
\[
\langle x \rangle \underset{\text{def}}{=} \{ e, x, \dots, x^{\operatorname{Ord}(x)-1} \}.
\]
\end{definition}

\begin{lemma}[{Cardinality of $\langle x \rangle$}]\label[lemma]{lem:cyclic-subgroup-order}
The elements listed in \cref{def:cyclic-subgroup} are all pairwise distinct, so
\[
|\langle x \rangle| = \bigl|\{ e, x, \dots, x^{\operatorname{Ord}(x)-1} \}\bigr| = \operatorname{Ord}(x).
\]
\end{lemma}

\begin{proof}
Write $m := \operatorname{Ord}(x)$ and $x^0 := e$. There are $m$ expressions $x^0, x^1, \dots, x^{m-1}$, so we must show they are pairwise distinct.

First, no $x^j$ with $1 \le j \le m - 1$ can equal $e$, because $m$ is by definition the minimal positive integer with $x^m = e$.

More generally, suppose $x^i = x^j$ for some $0 \le i < j \le m - 1$. By $i$-fold cancellation (\cref{prop:cancellation-laws}), we obtain
\[
e = x^{j - i}.
\]
Since $0 < j - i \le m - 1 < m$, this contradicts the minimality of $m = \operatorname{Ord}(x)$.

Hence, all $m$ elements $e, x, \dots, x^{\operatorname{Ord}(x)-1}$ are distinct, and $|\langle x \rangle| = \operatorname{Ord}(x)$.
\end{proof}
\end{content}

\begin{hidden}
% timestamp: 01:27:11
% topic: Defined <x> and proved |<x>| = ord(x)
% board_state: def:cyclic-subgroup, lem:cyclic-subgroup-order
% next_goal: State that <x> is a subgroup isomorphic to Z/ord(x)Z and close the lecture
% open_loops: full proof that <x> is a subgroup and <x> isomorphic to Z/ord(x)Z (postponed to next lecture)
\end{hidden}

\begin{speech}[01:27:11 - 01:27:30]
So, let me say the final thing, then the bell's going to ring.

Actually, this set is a subgroup, and moreover this subgroup is isomorphic to something \inlinemetanote{writes \qt{$\langle x \rangle \subset G \text{ is a subgroup}$} and \qt{$\langle x \rangle \cong$}}. And what is it isomorphic to? Yeah.
\end{speech}

\begin{student}[Student Answer]
$\mathbb{Z}$ mod $\operatorname{Ord}(x)$.
\end{student}

\begin{speech}[01:27:32 - 01:28:18]
$\mathbb{Z}$ mod $\operatorname{Ord}(x)\mathbb{Z}$ \inlinemetanote{writes: $\faktor{\mathbb{Z}}{\operatorname{Ord}(x)\mathbb{Z}}$}.

So, it's a subgroup: it says that if I multiply anything here or their inverse, I... I have to first show it's closed under multiplication, and it closes under inverse --- I'll do that next time. It's kind of basic, so I'll... I mean, it's kind of important, so I'll do it next time.

But the first thing is you have to show when you multiply these together, you stay in this set, and you have to show the inverses are there, and finally what the subgroup is is just this with addition and zero \inlinemetanote{places parentheses around the quotient and writes \qt{$, +, 0)$}, giving $(\faktor{\mathbb{Z}}{\operatorname{Ord}(x)\mathbb{Z}}, +, 0)$}.

So, I'll explain that next time.
\end{speech}

\begin{content}[{Subgroup Property and Isomorphism of $\langle x \rangle$}]
\begin{proposition}[{Cyclic Subgroup Generated by $x$}]\label[proposition]{prop:cyclic-isomorphism}
Let $G$ be a finite group and $x \in G$. Then:
\begin{enumerate}[label=\textbf{(\arabic*)}]
    \item $\langle x \rangle \subset G$ is a subgroup of $G$.
    \item As a group, $\langle x \rangle$ is isomorphic to the additive group of integers modulo $\operatorname{Ord}(x)$:
    \[
    \langle x \rangle \cong \left(\faktor{\mathbb{Z}}{\operatorname{Ord}(x)\mathbb{Z}}, +, 0\right).
    \]
\end{enumerate}
\end{proposition}

\begin{steps}[Roadmap for the Next Lecture]
The lecturer proves \cref{prop:cyclic-isomorphism} in the next lecture (\cref{prop:cyclic-subgroup}, \cref{cor:cyclic-subgroup-isomorphism}), checking three points:
\begin{itemize}
    \item \textbf{Closure under multiplication:} the product of any two elements in $\langle x \rangle$ stays in $\langle x \rangle$ (using division with remainder $i + j = qm + r$ for $m = \operatorname{Ord}(x)$).
    \item \textbf{Closure under inverses:} the inverse of $x^i \in \langle x \rangle$ ($0 < i < m$) is $x^{m-i} \in \langle x \rangle$.
    \item \textbf{Isomorphism:} showing that $\langle x \rangle \cong \bigl(\faktor{\mathbb{Z}}{\operatorname{Ord}(x)\mathbb{Z}}, +, 0\bigr)$ via the explicit isomorphism $\bar{k} \mapsto x^k$ (well-defined because $x^m = e$, and bijective because both sides have cardinality $m$ by \cref{lem:cyclic-subgroup-order}).
\end{itemize}
\end{steps}
\end{content}

\begin{aiNote}[Beyond the Lecture: Cyclic Subgroups]
% [PERSISTENT, ADDITION]
The subgroup $\langle x \rangle$ is also called the \emph{cyclic subgroup} generated by $x$, and it can equivalently be written as $\{ x^k \mid k \in \mathbb{Z} \}$.
\end{aiNote}

\begin{overview}[Summary of Key Concepts]
\begin{itemize}
    \item \textbf{Gauss's Theorem (Theorem~B) and Examples of UFDs:} If $R$ is a UFD, then $R[x]$ is a UFD (\cref{thm:ufd-polynomial-ring-w7tue}). Iterating yields UFDs such as $\mathbb{Z}[x], \mathbb{Z}[x, y], \dots$, $\mathbb{F}_p[x, y, z, \dots]$, and $\mathbb{Q}[x, y, z]$ (\cref{exm:classical-ufds}). In addition, $\mathbb{Z}[i]$ (\cref{ex:gaussian-integers-ufd}) and the formal power series ring $k[[x]]$ over a field $k$ (\cref{ex:power-series-pid-ufd}) are PIDs and hence UFDs; every non-zero ideal of $k[[x]]$ is of the form $(x^m)$ with $m \ge 0$ (\cref{exer:ideals-power-series}).
    \item \textbf{Content, Primitive Polynomials, and Gauss's Lemma:} Any non-zero $f \in R[x]$ decomposes as $f = C(f) f^*$ where $C(f)$ is the $\operatorname{GCD}$ of its non-zero coefficients up to a unit and $C(f^*) \stackrel{u}{=} 1$ (\cref{def:content-primitive}). Products of primitive polynomials are primitive (\cref{prop:gauss-lemma-product}), and if a primitive $g \in R[x]$ divides $bf$ ($0 \neq b \in R$), then $g \mid f$ (\cref{prop:primitive-divisibility}).
    \item \textbf{Proof of Case (ii) ($\deg(p) > 0$) of Gauss's Theorem:} Given $p \in R[x]$ irreducible of positive degree, $fg \in (p)$, and $f \notin (p)$, choosing $m \in I = (f, p)$ of minimal degree, dividing $f$ and $p$ by $m$ in $Q(R)[x]$, clearing denominators, and using irreducibility of $p$ forces $m \in R \setminus \{0\}$. Writing $mg = Afg + Bgp$ gives $p \mid mg$; since $\deg(p) > 0$ implies $p$ is primitive (\cref{clm:irr-pos-deg-primitive}), it follows that $p \mid g$, so $g \in (p)$.
    \item \textbf{Group Axioms and Elementary Properties:} A group $(G, \cdot, e)$ satisfies associativity \textbf{(G1)}, the unit property \textbf{(G2)}, and existence of inverses \textbf{(G3)} (\cref{def:group-w7tue}). The unit and inverses are unique (\cref{prop:group-unit-unique} and \cref{prop:group-inverse-unique}), left and right cancellation hold (\cref{prop:cancellation-laws}), and counting parenthesizations of an $n$-fold product gives the Catalan number $C_{n-1}$, where $C_n = \frac{1}{n+1}\binom{2n}{n}$ (\cref{ans:catalan-brackets}).
    \item \textbf{Subgroups, Homomorphisms, and $S_n$:} A subgroup $S \subset G$ is closed under multiplication, contains $e$, and is closed under $\operatorname{Inv}$ (\cref{def:subgroup-w7tue}). Any group homomorphism $f \colon G_1 \to G_2$ (\cref{def:group-homomorphism-w7tue}) satisfies $f(e_1) = e_2$ and $f(x^{-1}) = f(x)^{-1}$ (\cref{prop:hom-preserves-structure}). The symmetric group $S_n$ of bijections $[n] \to [n]$ has order $n!$ (\cref{def:symmetric-group} and \cref{prop:order-symmetric-group}); cycle notation (\cref{def:cycle}, \cref{not:omitting-1-cycles}) efficiently represents its elements and illustrates that $S_3$ is non-abelian (\cref{exm:s3-cycles}, \cref{obs:s3-non-commutative}).
    \item \textbf{Orders of Elements in Finite Groups:} For every element $x$ in a finite group $G$, there exists $m > 0$ with $x^m = e$ (\cref{prop:finite-group-identity-power}), and the minimal such $m$ is $\operatorname{Ord}(x)$ (\cref{def:element-order}). The set $\langle x \rangle = \{e, x, \dots, x^{\operatorname{Ord}(x)-1}\}$ has cardinality $\operatorname{Ord}(x)$ (\cref{def:cyclic-subgroup} and \cref{lem:cyclic-subgroup-order}) and forms a subgroup of $G$ isomorphic to $(\faktor{\mathbb{Z}}{\operatorname{Ord}(x)\mathbb{Z}}, +, 0)$ (\cref{prop:cyclic-isomorphism}).
\end{itemize}
\end{overview}